\documentclass[
12pt,
fancyheadings,
a4paper,
twoside,
BCOR=0.7cm,
nobcorignoretitle,
grayscalebody,
]{stsreprt}

\author{Marc Bessa Hoffmann}
\title{Efficient and Interpretable Dynamic Slicing of Neural Networks: A Block-Based Approach}
\date{\today}
\subject{Master Thesis}
\professor{Prof. Dr. Sibylle Schupp}
\advisor{Daniel Rashedi}

\usepackage[utf8]{inputenc}
\usepackage[T1]{fontenc}
\usepackage{lmodern}
\usepackage{listings}
\usepackage{xcolor}
\usepackage{amsmath}
\usepackage{amssymb}
\usepackage{amsthm}
\usepackage{booktabs}
\usepackage{multirow}
\usepackage{floatbytocbasic}
\usepackage{algorithm}
\usepackage{algpseudocode}
\usepackage{tikz}
\usepackage{wrapfig}
\usepackage{pgfplots}
\usepackage{url}
\pgfplotsset{compat=1.18}
\usepackage{lipsum}
\usepackage[
    backend=biber,
    style=numeric,
    sorting=none,
    maxnames=100,
    minnames=1,
    doi=true,
    url=true,
    eprint=false
]{biblatex}
\addbibresource{references.bib}

\definecolor{staticslice}{RGB}{214, 234, 248}
\definecolor{dynamicslice}{RGB}{255, 235, 205}
\definecolor{codegreen}{rgb}{0,0.6,0}
\definecolor{codegray}{rgb}{0.5,0.5,0.5}
\definecolor{codepurple}{rgb}{0.58,0,0.82}
\definecolor{backcolour}{rgb}{0.95,0.95,0.92}

\lstdefinestyle{python}{
        language=Python,
        backgroundcolor=\color{backcolour},
        commentstyle=\color{codegreen},
        keywordstyle=\color{magenta},
        numberstyle=\tiny\color{codegray},
        stringstyle=\color{codepurple},
        basicstyle=\ttfamily\small,
        breakatwhitespace=false,
        breaklines=true,
        captionpos=b,
        keepspaces=true,
        numbers=left,
        numbersep=5pt,
        showspaces=false,
        showstringspaces=false,
        showtabs=false,
        tabsize=4,
}
\newtheorem{definition}{Definition}[chapter]

% Increase vertical space around equations
\setlength{\abovedisplayskip}{14pt}
\setlength{\belowdisplayskip}{14pt}
\setlength{\abovedisplayshortskip}{8pt}
\setlength{\belowdisplayshortskip}{8pt}

\begin{document}

\frontmatter
\maketitle

\newpage\null\thispagestyle{empty}\newpage

\begin{center}
\section*{Statutory Declaration}
\end{center}

\vspace{1cm}

\noindent I herewith declare that I have composed the present thesis myself and without use of any other than the
cited sources and aids. Sentences or parts of sentences quoted literally are marked as such; other
references
with regard to the statement and scope are indicated by full details of the publications concerned.
The thesis
in the same or similar form has not been submitted to any examination body and has not been
published.
This thesis was not yet, even in part, used in another examination or as a course performance.

\vspace{2cm}

\noindent
\begin{tabular}{@{}p{0.48\textwidth}p{0.48\textwidth}@{}}
Hamburg, March 30, 2026 &  \_\_\_\_\_\_\_\_\_\_\_\_\_\_\_\_ \\
& Marc Bessa Hoffmann
\end{tabular}

\vfill

\begin{center}
\footnotesize Text source: Goethe University, Department of Economics
\end{center}

\chapter*{\centering \begin{normalsize}Abstract\end{normalsize}}
\begin{quotation}
Neural networks achieve strong performance in many domains but are often considered difficult to
interpret.
Understanding which internal components contribute to a specific prediction is an active research
topic in explainable artificial intelligence (XAI).
Program slicing has recently been adapted from traditional software systems to address this
problem.
NNSlicer propagates contribution scores for each neuron backward from the prediction output through
the network, yielding a slice: a subset of neurons and synapses that are most relevant to a given
prediction.
Despite its success, NNSlicer operates at a fine-grained neuron level and is time-consuming, taking
several minutes to slice even smaller ResNet models.

To improve runtime efficiency, we build on the backward contribution analysis of NNSlicer and
propose a hierarchical block-based slicing approach that leverages the structural layout of
residual networks.
For each input, we measure the deviation of neuron responses from their dataset averages.
We aggregate these deviations into channel-level and block-level indicators that express how
strongly each component reacts to the input.
Based on these indicators, we first filter inactive blocks, then inactive channels, and only then
perform neuron-level analysis on the remaining components.
For the neuron-level analysis, we propose a vectorized implementation that operates on tensors to
further improve runtime.

Our vectorized implementation alone slices ResNet-18 in about two seconds compared to over 540
seconds reported for NNSlicer.
Even ResNet-101, which was previously infeasible to slice in practice, can now be processed in
about 7.7 seconds.
Block-based slicing further increases efficiency while preserving slice quality: for ResNet-101,
runtime drops from 7.7\,s to 6.2\,s and the slice size decreases from 41.8\,\% to 23.5\,\%.
Slice quality is evaluated on a custom ResNet with 27 residual blocks and 855K parameters.
Model pruning shows that a block-pruned model retains 91.9\% target accuracy with only 39.7\% of
its parameters.
Counterfactual experiments further confirm that progressively removing neurons inside the slice
degrades the prediction, whereas the prediction remains stable when removing neurons outside the
slice.
Finally, we observe that channels and blocks in the last residual blocks exhibit interpretable
patterns across slices of semantically related classes.
\end{quotation}
\tableofcontents
\listoffigures{}
\listoftables{}
\lstlistoflistings{}
\listofalgorithms{}
\mainmatter


\chapter{Introduction}

Neural networks are widely used in domains such as medical imaging, autonomous driving, and natural
language processing~\cite{lecun2015deep}.
Despite their strong performance, they are often treated as black boxes whose internal decision
process is difficult to understand.
In safety-critical settings, the lack of transparency makes it hard to assess the reliability of
individual predictions and highlights the need for trustworthy models.
This need has motivated the field of explainable artificial intelligence (XAI), which aims to make
model decisions more understandable to humans~\cite{Samek_2021}.

In recent years, program slicing has emerged as an approach within XAI that focuses on internal
components of neural networks.
Program slicing is a technique that was originally developed for traditional software systems.
The core idea is to identify the parts of a program that are relevant to a particular computation
or output~\cite{weiser1984program}.
Over several decades, it has been successfully applied to tasks such as debugging, program
comprehension, and testing~\cite{10.1145/1050849.1050865}.
The concept of program slicing has been adapted to neural networks to analyze which internal
components contribute most to a specific
prediction~\cite{wang2018interpret,zhang2020dynamic,zhou2024neusemslice}.

Among these, NNSlicer is a dynamic program slicing approach that operates in three
phases~\cite{zhang2020dynamic}.
First, a profiling phase runs the training data through the network and computes the mean
activation at each neuron.
In the second phase, forward analysis compares the activations for a specific input against these
means. The resulting delta for each neuron indicates how much its activation deviates from the
average.
Finally, the backward analysis traces contributions from the output back to identify the relevant
neurons and synapses.
The output of NNSlicer is a slice of a neural network that identifies the most relevant neurons and
synapses with positive or negative contributions.
NNSlicer shows the utility of such slices for applications like adversarial input detection,
targeted model pruning, and selective model protection.

Although NNSlicer demonstrates that program slicing can be applied to neural networks, it operates
on a fine-granular neuron level.
While this provides precise per-neuron contributions, modern architectures can contain millions of
parameters distributed across hundreds of layers.
Slicing such models at the neuron level is time-consuming, taking several minutes even for smaller
models and becoming infeasible for larger ones.
Making slicing more efficient in terms of runtime is therefore our primary objective.

In this thesis, we build on NNSlicer and propose a hierarchical block-based slicing approach that
leverages the structural layout of the neural network.
We focus exclusively on convolutional residual networks (ResNets), as they provide a clear
residual-block structure for higher-level slicing.
Our method first filters at the residual block level, then at the channel level, and finally
performs neuron-level analysis on the remaining components.
To determine which blocks and channels to retain, we extend the forward analysis of NNSlicer to
aggregate neuron deltas not only per neuron, but also per channel and per residual block.
These aggregated deltas indicate how strongly each component reacts to a given input compared to
its average behavior.
Based on these aggregated deltas, we rank blocks and channels in descending order and retain only
those that account for a configurable fraction of the total delta. Blocks and channels below this
threshold are filtered out and excluded from the slice.
For the blocks and channels that are included in the slice, we compute the neuron-level
contributions within each block and channel as NNSlicer does. 

Since the original NNSlicer iterates over individual neurons in Python loops, the computation time
grows with the number of neurons and does not scale to larger models. 
As a second contribution alongside block-based slicing, we propose a vectorized PyTorch
reimplementation that computes neuron contributions through tensor operations instead of processing
them one by one.
Together, block-based slicing and the vectorized implementation aim to increase runtime efficiency
and make neural network slicing practical for larger models.

Our evaluation is structured around three research questions.
The first research question focuses on runtime performance and slice size.
To assess slice size, we measure the proportion of active neurons and synapses within a slice
compared to the entire network.
We evaluate our method on ResNet-18, ResNet-34, and ResNet-101 to cover models of increasing depth
and complexity.
Our vectorized reimplementation slices a ResNet-18 in approximately two seconds, while NNSlicer
reports over 540 seconds for the same model.
Even larger models such as ResNet-101 with 44.5 million parameters, previously infeasible to slice
with NNSlicer, can now be sliced in about 7.7 seconds with our vectorized implementation.
Block-based slicing further reduces both runtime and slice size.
For a ResNet-101 slice, we achieve a runtime improvement of 19\,\% by dropping from 7.67\,s to
6.20\,s.
Block-based slicing is particularly effective at producing smaller slices. For the same ResNet-101
slice, the slice size reduces from 41.8\,\% without block-based slicing to 23.5\,\% with
block-based slicing.


Since efficiency alone is not sufficient, the second research question examines whether block-based
slicing preserves slice quality.
By quality, we mean whether the slice still captures the neurons and synapses that are most
relevant to the prediction.
The evaluation of slice quality is conducted on a custom ResNet that has been trained on the
CIFAR-10 dataset and consists of 27 residual blocks with 855K parameters.
The architecture follows the CIFAR-10 ResNet design described by He et al.~\cite{he2016deep}, with
three groups of nine BasicBlocks each and only identity shortcuts.
We chose a model with 27 residual blocks to provide sufficient depth for evaluating the block-based
slicing approach.
CIFAR-10 was chosen because it is a widely used benchmark that was also used in the original
NNSlicer evaluation.
It can be automatically downloaded with PyTorch and makes our results fully reproducible.

We use two different experiments to assess quality.
In model pruning, we remove all components outside the slice and measure whether the pruned model
still correctly classifies the target class.
Block-based slicing retains 91.9\,\% target accuracy while reducing the model to 39.7\,\% of its
original size.
As a second experiment, we use counterfactual reasoning, which asks what would happen if parts of
the network were different.
We examine two aspects: necessity tests whether removing neurons inside the slice degrades the
prediction, and sufficiency tests whether the prediction remains stable when neurons outside the
slice are removed.
Block-based slicing yields counterfactual behavior nearly identical to slicing without any
filtering while producing smaller slices. In contrast, aggressive block-based slicing shows a rapid
drop in prediction confidence, confirming that the filtered components are indeed relevant.

The third research question investigates interpretability at the block and channel level.
As an exploratory step, we examine whether patterns can be observed between slices of different
classes.
We compare classes where similarities are expected: automobile and truck as vehicles, and cat and
dog as animals.
On the block level, we observe overlaps in the final blocks of the network.
When zooming into the last block, we find channels with strongly shared contributions between cat
and dog as well as between automobile and truck.

Overall, our results show that block-based slicing and our vectorized reimplementation of NNSlicer
substantially reduce runtime and slice size while preserving slice quality.
We also find that the resulting slices support interpretability by highlighting class-specific
patterns in the last residual blocks and their most influential channels.

In summary, the contributions of this thesis are as follows:
\begin{itemize}
  \item Block-based slicing approach: a hierarchical extension that filters at the residual block and channel level before performing neuron-level analysis.
  \item Vectorized PyTorch implementation of NNSlicer that replaces neuron-by-neuron loops with tensor operations.
  \item Evaluation of runtime and slice size across ResNet models of increasing depth. Slice quality and interpretability are evaluated on a custom ResNet model using model pruning and counterfactual reasoning.
\end{itemize}

\vspace{0.2cm}
We structure this thesis as follows.
First, we introduce neural networks and the architectures used in this work, namely convolutional
and residual networks.
We then describe program slicing and how the concept can be applied to neural networks.
Next, we review related work. Since our approach builds directly on NNSlicer, we dedicate a
separate chapter to its pipeline.
We then introduce our block-based slicing approach and the vectorized implementation.
Finally, we present the evaluation results and conclude with a summary of findings and directions
for future work.



\chapter{Background}
\label{chap:background}

In this chapter, we provide the technical background and definitions used throughout this work.
We start by introducing the concept of a neuron and how neurons are combined into a neural network.
After that, we describe two neural network architectures that are of particular interest:
convolutional neural networks and residual networks.
Finally, we discuss program slicing and how it can be applied to neural networks.

\section{Neural Networks}
\label{sec:neural_networks}

Neural networks are mathematical models that draw inspiration from biological neural systems.
In the brain, a neuron fires when the sum of its incoming signals exceeds a threshold.
The same principle underlies the mathematical neuron, where a linear combination is calculated for
a set of numerical inputs to produce an output~\cite{rosenblatt1958perceptron}.
If we consider a neuron with $d$ inputs, the input vector is defined as:
\begin{equation}
    \mathbf{x} = (x_1, x_2, \dots, x_d)^\top
\end{equation}
Each input $x_i$ has an associated weight $w_i$:
\begin{equation}
    \mathbf{w} = (w_1, w_2, \dots, w_d)^\top
\end{equation}
The neuron computes a weighted sum of all inputs and adds a bias term $b$:
\begin{equation}
        z = \mathbf{w}^\top \mathbf{x} + b = \sum_{i=1}^{d} w_i x_i + b
\end{equation}
The weights determine how much influence each input has on the output.
For example, a neuron that decides whether to bring an umbrella might assign a high weight to the
input \textit{is it raining} and a low weight to \textit{is it sunny}.
While the weights control the influence of each input, the bias $b$ shifts the activation threshold
independently of the inputs.
If we consider the umbrella example again, a negative bias would mean that the evidence for rain
must be strong before the neuron decides to bring one.
A positive bias makes the neuron more cautious, activating even with a weak signal for \textit{is
it raining}.

As defined above, $z$ is so far only a linear combination of inputs and can only represent linear
transformations.
To solve more complex tasks, non-linearity is introduced by passing $z$ through an
\textit{activation function} $\sigma$.
The output of the neuron is therefore:
\begin{equation}
        y = \sigma(z)
\end{equation}
Many different activation functions exist for different purposes.
For instance, the sigmoid activation function maps values between the range of 0 and 1.
The output is therefore bounded to a fixed range, which prevents activations from growing across
layers.
However, such functions can make training in deep networks difficult.
Since the output is bounded to a small range, the gradients during backpropagation can become very
small.
As these small gradients are multiplied across many layers, they shrink further and can effectively
vanish~\cite{Goodfellow-et-al-2016}.
Rectified linear functions address this problem and have become the standard choice in modern
architectures~\cite{krizhevsky2017imagenet}.
A survey covering further activation functions is provided by Dubey et
al.~\cite{dubey2022activation}.

In this work, all models that have been evaluated use the ReLU activation function (Rectified
Linear Unit)~\cite{nair2010rectified}:
\begin{equation}
    \text{ReLU}(z) = \max(0, z) = \begin{cases} z, & \text{if } z \geqslant 0 \\ 0, & \text{otherwise} \end{cases}
\end{equation}
ReLU is applied to the raw output $z$ of the linear combination of a neuron: values where $z
\geqslant 0$ are passed through for further processing, while negative values are set to zero.
Besides ReLU, another important activation function is Softmax~\cite{Goodfellow-et-al-2016}:
\begin{equation}
    \text{Softmax}(z_i) = \frac{e^{z_i}}{\sum_{j=1}^{K} e^{z_j}}
\end{equation}

\noindent
The Softmax activation function is used for the prediction in the final output layer. It takes the
raw values of all output neurons and normalizes them into a probability distribution.
For example, given three output neurons with raw values $(2.0, 1.0, 0.1)$, Softmax produces
approximately $(0.66, 0.24, 0.10)$, indicating a 66\% probability for the first class.

With activation functions in hand, we can now formalize the definitions used in this work.
The following definitions for neural networks are based on the concepts introduced by Goodfellow et
al.~\cite{Goodfellow-et-al-2016} and have been adapted for this work. The notation is aligned with
NNSlicer, as our work builds directly on their formulation~\cite{zhang2020dynamic}.

\begin{definition}[Neuron]
\label{def:neuron}
A neuron $n$ takes an input vector $\mathbf{x} \in \mathbb{R}^d$ and a weight vector $\mathbf{w}
\in \mathbb{R}^d$ and computes the output $y = \sigma(z)$, where $z = \mathbf{w}^\top \mathbf{x} +
b$ is the weighted sum with bias $b \in \mathbb{R}$, and $\sigma$ is a non-linear activation
function.
We denote the set of all neurons in a network as $\mathcal{N}$.
\end{definition}

In practice, a single neuron is not used in isolation.
Multiple neurons are connected to form a neural network, where the output of one neuron serves as
input to another.
These connections carry the weights introduced above and are called \textit{synapses}.

\begin{definition}[Synapse]
\label{def:synapse}
A synapse $s = (n_i, n_j, w)$ is a directed weighted connection from neuron~$n_i$ to neuron~$n_j$
with weight $w \in \mathbb{R}$.
The value transmitted through the synapse is the activation $y_i$ of the source neuron~$n_i$.
We denote the set of all synapses in a network as $\mathcal{S}$.
\end{definition}

Neurons and synapses in a network are organized into layers.
The first layer is the input layer, which receives the raw input data.
The last layer is the output layer, which produces the final prediction.
Layers between input and output are called hidden layers, as their activations are internal to the
network and not directly observable.

\begin{definition}[Layer]
\label{def:layer}
A layer $l = (\mathcal{N}_l, \mathcal{S}_l)$ consists of a subset of neurons $\mathcal{N}_l \subset
\mathcal{N}$ and a subset of synapses $\mathcal{S}_l \subset \mathcal{S}$. The synapses
$\mathcal{S}_l$ are all synapses whose target neuron belongs to $\mathcal{N}_l$.
We denote the set of all layers in a network as $\mathcal{L}$.
\end{definition}

Layers differ in how their neurons are connected to the preceding layer.
In a fully connected layer, every neuron receives input from all neurons in the preceding layer.

\begin{definition}[Fully Connected Layer]
\label{def:fc_layer}
A fully connected layer $l = (\mathcal{N}_l, \mathcal{S}_l)$ is a layer in which every neuron in
$\mathcal{N}_l$ is connected to every neuron in the preceding layer $\mathcal{N}_{l-1}$, i.e., for
every pair $(n_i, n_j)$ with $n_i \in \mathcal{N}_{l-1}$ and $n_j \in \mathcal{N}_l$, there exists
a synapse $(n_i, n_j, w) \in \mathcal{S}_l$.
\end{definition}

Figure~\ref{fig:neural_network} shows an example of a simple neural network with fully connected
layers. It consists of $|\mathcal{L}| = 4$ layers, $|\mathcal{N}| = 8$ neurons, and $|\mathcal{S}|
= 14$ synapses. The input layer contains two neurons (blue), followed by two hidden layers with 3
and 2 neurons (yellow). The output layer has one neuron (green). The gray edges represent the
synapses.
\begin{figure}[ht]
    \centering
    \begin{tikzpicture}[scale=1.1, every node/.style={font=\footnotesize},
                neuron/.style={circle, draw, minimum size=1cm, inner sep=0pt}]
      \node[neuron, fill=blue!10] (n11) at (0, 0.9) {$1{,}1$};
      \node[neuron, fill=blue!10] (n12) at (0, -0.9) {$1{,}2$};
      \node[neuron, fill=yellow!20] (n21) at (2.5, 1.8) {$2{,}1$};
      \node[neuron, fill=yellow!20] (n22) at (2.5, 0) {$2{,}2$};
      \node[neuron, fill=yellow!20] (n23) at (2.5, -1.8) {$2{,}3$};
      \node[neuron, fill=yellow!20] (n31) at (5, 0.9) {$3{,}1$};
      \node[neuron, fill=yellow!20] (n32) at (5, -0.9) {$3{,}2$};
      \node[neuron, fill=green!15] (n41) at (7.5, 0) {$4{,}1$};
      \draw[gray!70] (n11) -- (n21);
      \draw[gray!70] (n11) -- (n22);
      \draw[gray!70] (n11) -- (n23);
      \draw[gray!70] (n12) -- (n21);
      \draw[gray!70] (n12) -- (n22);
      \draw[gray!70] (n12) -- (n23);
      \draw[gray!70] (n21) -- (n31);
      \draw[gray!70] (n21) -- (n32);
      \draw[gray!70] (n22) -- (n31);
      \draw[gray!70] (n22) -- (n32);
      \draw[gray!70] (n23) -- (n31);
      \draw[gray!70] (n23) -- (n32);
      \draw[gray!70] (n31) -- (n41);
      \draw[gray!70] (n32) -- (n41);
      \node[below, font=\small] at (0, -2.4) {Input};
      \node[below, font=\small] at (2.5, -2.4) {Hidden};
      \node[below, font=\small] at (5, -2.4) {Hidden};
      \node[below, font=\small] at (7.5, -2.4) {Output};
    \end{tikzpicture}
    \caption{A neural network with fully connected layers. An input layer with two neurons. Two hidden layers with three and two neurons. One output neuron for the prediction. The edges represent the synapses.}
    \label{fig:neural_network}
\end{figure}

More formally, a neural network can be described by its layers, neurons, and the synapses
connecting them.

\begin{definition}[Neural Network]
\label{def:neural_network}
A neural network is a tuple $\mathcal{M} = (\mathcal{N}, \mathcal{S})$, where $\mathcal{N}$ is the
set of neurons and $\mathcal{S} \subseteq \mathcal{N} \times \mathcal{N} \times \mathbb{R}$ is the
set of synapses.
The neurons are organized into layers $\mathcal{L} = \{l_1, \dots, l_L\}$: an input layer, one or
more hidden layers, and an output layer.
\end{definition}

The input to a neural network $\mathcal{M}$ is a sample $\xi$, a numerical vector representing a
single data point. We define a dataset as a collection of such samples $\mathcal{I} = \{\xi_1,
\xi_2, \dots, \xi_m\}$.

The weights of a neural network must be trained before the model can be used for predictions. The
following paragraph gives an overview of the training process.

\paragraph{Training.}
Initially, the weights of a network are set to small random values.
The network learns by adjusting its weights and biases to minimize a loss function that measures
the difference between the predicted output and the expected output.
For classification tasks, a common choice is the cross-entropy loss.
Given a true class label $y$ and the predicted probability $p_y$ from the Softmax output, it is
defined as:
\begin{equation}
        L = -\log(p_y)
\end{equation}
The loss is small when the predicted probability for the correct class is high, and large when it
is low.
To minimize the loss, the network needs to determine how much each weight contributes to the error.
This is computed using backpropagation, which applies the chain rule of calculus to propagate the
gradient of the loss backwards through the network, layer by layer~\cite{rumelhart1986learning}.
For a loss $L$ and a weight $w$ in layer $l$, the chain rule decomposes the gradient as:
\begin{equation}
    \frac{\partial L}{\partial w} = \frac{\partial L}{\partial z_l} \cdot \frac{\partial z_l}{\partial w}
\end{equation}
where $z_l$ is the output of layer $l$.
Once the gradients are computed, the weights are updated using gradient descent.
In each step, a weight $w$ is updated by subtracting the gradient scaled by a learning rate $\eta$:
\begin{equation}
        w \leftarrow w - \eta \cdot \frac{\partial L}{\partial w}
\end{equation}
The learning rate $\eta$ controls how large each update step is. A value too large can cause the
training to overshoot, while a value too small makes training unnecessarily slow.
Several optimizers have been developed to improve gradient descent, such as SGD with momentum or
Adam~\cite{qian1999momentum, kingma2015adam}.
During training with such an optimizer, the network processes the dataset over multiple epochs. One
epoch means every sample $\xi \in \mathcal{I}$ has been processed once.
In practice, $\mathcal{I}$ is split into smaller batches, and the weights are updated after each
batch.

Once the training is complete, the network can make predictions by passing new inputs through the
network in a single forward pass.
The slicing approach presented in this work operates on this forward pass, using the intermediate
activation values to determine which parts of the network are relevant for a given prediction.

The neural networks described so far use fully connected layers, where each neuron in one layer is
connected to every neuron in the next.
The following section introduces convolutional neural networks, an alternative architecture
specialized for processing images.

\section{Convolutional Neural Networks}

Convolutional neural networks (CNNs) are a class of neural networks designed for processing
structured data such as images~\cite{lecun1998gradient}.
The core idea is the convolution operation: a small matrix called a filter (or kernel) with weights
$\mathbf{w}$ and bias $b$ slides over the input and computes a weighted sum at each position.
Each filter is designed to detect a specific pattern in the input.
For example, a filter can detect vertical edges by responding strongly wherever the pixel value
changes from left to right.
The convolution operation is illustrated in Figure~\ref{fig:convolution}.

\begin{figure}[ht!]
    \centering
    \begin{tikzpicture}[scale=0.8, every node/.style={font=\footnotesize}]
    \foreach \r in {0,...,4} {
      \foreach \c in {0,...,4} {
        \draw[gray!60] (\c, -\r) rectangle ++(1,1);
            }
        }
    \foreach \r in {0,...,2} {
      \foreach \c in {0,...,2} {
        \fill[blue!20] (\c, -\r) rectangle ++(1,1);
        \draw[blue!70, thick] (\c, -\r) rectangle ++(1,1);
            }
        }
    \node at (0.5, 0.5) {1}; \node at (1.5, 0.5) {0}; \node at (2.5, 0.5) {1}; \node at (3.5, 0.5) {0}; \node at (4.5, 0.5) {2};
    \node at (0.5,-0.5) {1}; \node at (1.5,-0.5) {1}; \node at (2.5,-0.5) {0}; \node at (3.5,-0.5) {1}; \node at (4.5,-0.5) {0};
    \node at (0.5,-1.5) {0}; \node at (1.5,-1.5) {1}; \node at (2.5,-1.5) {1}; \node at (3.5,-1.5) {0}; \node at (4.5,-1.5) {1};
    \node at (0.5,-2.5) {2}; \node at (1.5,-2.5) {0}; \node at (2.5,-2.5) {1}; \node at (3.5,-2.5) {1}; \node at (4.5,-2.5) {0};
    \node at (0.5,-3.5) {0}; \node at (1.5,-3.5) {1}; \node at (2.5,-3.5) {0}; \node at (3.5,-3.5) {2}; \node at (4.5,-3.5) {1};
    \node[above] at (2.5, 1) {\footnotesize Input};
    \begin{scope}[xshift=7.5cm, yshift=-1cm]
    \foreach \r in {0,...,2} {
      \foreach \c in {0,...,2} {
        \fill[orange!25] (\c, -\r) rectangle ++(1,1);
        \draw[orange!70, thick] (\c, -\r) rectangle ++(1,1);
            }
        }
    \node at (0.5, 0.5) {1}; \node at (1.5, 0.5) {0}; \node at (2.5, 0.5) {1};
    \node at (0.5,-0.5) {0}; \node at (1.5,-0.5) {1}; \node at (2.5,-0.5) {0};
    \node at (0.5,-1.5) {1}; \node at (1.5,-1.5) {0}; \node at (2.5,-1.5) {1};
    \node[above] at (1.5, 1) {\footnotesize Kernel};
    \end{scope}
    \draw[->, thick, gray] (5.3, -1.5) -- node[above] {\footnotesize $*$} (7.3, -1.5);
    \begin{scope}[xshift=13cm, yshift=-1cm]
    \foreach \r in {0,...,2} {
      \foreach \c in {0,...,2} {
        \draw[gray!60] (\c, -\r) rectangle ++(1,1);
            }
        }
    \fill[green!20] (0, 0) rectangle ++(1,1);
    \draw[green!70, thick] (0, 0) rectangle ++(1,1);
    \node at (0.5, 0.5) {4};
    \node at (1.5, 0.5) {2}; \node at (2.5, 0.5) {3};
    \node at (0.5,-0.5) {3}; \node at (1.5,-0.5) {3}; \node at (2.5,-0.5) {1};
    \node at (0.5,-1.5) {3}; \node at (1.5,-1.5) {3}; \node at (2.5,-1.5) {4};
    \node[above] at (1.5, 1) {\footnotesize Feature Map};
    \end{scope}
    \draw[->, thick, gray] (10.7, -1.5) -- node[above] {\footnotesize $=$} (12.8, -1.5);
    \end{tikzpicture}
    \caption{Convolution operation. A $3 \times 3$ kernel slides over a $5 \times 5$ input with stride~1 and no padding to produce a $3 \times 3$ feature map. The highlighted positions show the first computation}
    \label{fig:convolution}
\end{figure}

Stride and padding are two parameters that control the convolution and determine the dimensions of
the resulting feature map.
The stride defines how far the filter moves between positions.
Padding adds zeros around the input borders, which allows the output to preserve the spatial
dimensions of the input.
In the example above, a stride of 1 was used to move the filter by one position. No padding is
used.
The resulting feature map is $3 \times 3$ and is also known as a channel.
The terms \textit{feature map} and \textit{channel} are often used interchangeably. A feature map
refers to the spatial output of a convolution, while a channel refers to the structural unit within
a layer.

Similar to the neural network definitions, the following definitions are based on Goodfellow et
al.~\cite{Goodfellow-et-al-2016} and have been adapted for this work.

\begin{definition}[Channel]
\label{def:channel}
A channel $c$ is produced by applying a single filter with kernel weights $\mathbf{w} \in
\mathbb{R}^{k \times k}$ and bias $b \in \mathbb{R}$ to an input $\mathbf{x}$.
It consists of a set of neurons $\mathcal{N}_{ch}$ arranged in a spatial grid of height $H$ and
width $W$, connected by synapses $\mathcal{S}_{ch}$.
Each neuron $n_{i,j}$ computes a linear combination of a local $k \times k$ region of the input
$\mathbf{x}$ at position $(i, j)$.
\end{definition}

The feature map in Figure~\ref{fig:convolution} has $H = W = 3$, so the channel contains $9$
neurons in $\mathcal{N}_{ch}$.
Each neuron $n_{i,j}$ is connected to a $3 \times 3$ input region through $9$ synapses in
$\mathcal{S}_{ch}$, resulting in $9 \times 9 = 81$ synapses in total.
All $9$ neurons share the same kernel weights $\mathbf{w}$ and bias $b$.
Multiple filters are combined into a convolutional layer, where each filter produces one channel.

\begin{definition}[Convolutional Layer]
\label{def:conv_layer}
A convolutional layer $l$ consists of a set of channels $\mathcal{C}$, where each channel $c \in
\mathcal{C}$ is produced by a distinct filter.
The neurons of the layer are the union of all channel neurons $\mathcal{N}_l = \bigcup_{c \in
\mathcal{C}} \mathcal{N}_{ch}$, and similarly for the synapses $\mathcal{S}_l = \bigcup_{c \in
\mathcal{C}} \mathcal{S}_{ch}$.
\end{definition}

The convolutional layer in Figure~\ref{fig:cnn_architecture} has $|\mathcal{C}| = 3$ channels.
Each channel contains $30 \times 30 = 900$ neurons.
The layer therefore has $|\mathcal{N}_l| = 3 \times 900 = 2{,}700$ neurons in total.
With a $3 \times 3$ kernel, each neuron has $9$ synapses, giving $|\mathcal{S}_l| = 2{,}700 \times
9 = 24{,}300$ synapses for the entire layer.

Convolutional layers are typically combined with other layer types into a full architecture.
Figure~\ref{fig:cnn_architecture} shows an example: the input image is first passed to a
convolutional layer that produces three channels.
These are then reduced in size by a pooling layer.
A batch normalization layer normalizes the values before they are flattened and passed through
fully connected layers to produce the final classification.

\begin{figure}[ht!]
    \centering
    \begin{tikzpicture}[scale=0.7, every node/.style={font=\footnotesize}]
      \fill[gray!15] (0,0) rectangle (2.4,2.4);
      \draw[thick] (0,0) rectangle (2.4,2.4);
      \node at (1.2, -0.5) {Input};
      \node at (1.2, -1.0) {$32 \times 32 \times 1$};
      \draw[->, thick] (2.7, 1.2) -- (3.7, 1.2);
      \foreach \i in {0,1,2} {
        \pgfmathsetmacro{\xoff}{4.0 + \i*0.3}
        \pgfmathsetmacro{\yoff}{\i*0.3}
        \fill[blue!15] (\xoff, \yoff) rectangle ++(2.0, 2.0);
        \draw[blue!60, thick] (\xoff, \yoff) rectangle ++(2.0, 2.0);
            }
      \node at (5.3, -0.5) {Conv};
      \node at (5.3, -1.0) {$30 \times 30 \times 3$};
      \draw[->, thick] (6.9, 1.2) -- (7.9, 1.2);
      \foreach \i in {0,1,2} {
        \pgfmathsetmacro{\xoff}{8.2 + \i*0.3}
        \pgfmathsetmacro{\yoff}{0.35 + \i*0.3}
        \fill[red!15] (\xoff, \yoff) rectangle ++(1.3, 1.3);
        \draw[red!60, thick] (\xoff, \yoff) rectangle ++(1.3, 1.3);
            }
      \node at (9.2, -0.5) {Pool};
      \node at (9.2, -1.0) {$15 \times 15 \times 3$};
      \draw[->, thick] (10.4, 1.2) -- (11.4, 1.2);
      \foreach \i in {0,1,2} {
        \pgfmathsetmacro{\xoff}{11.7 + \i*0.3}
        \pgfmathsetmacro{\yoff}{0.35 + \i*0.3}
        \fill[orange!20] (\xoff, \yoff) rectangle ++(1.3, 1.3);
        \draw[orange!60, thick] (\xoff, \yoff) rectangle ++(1.3, 1.3);
            }
      \node at (12.7, -0.5) {Batch Norm};
      \node at (12.7, -1.0) {$15 \times 15 \times 3$};
      \draw[->, thick] (13.9, 1.2) -- (14.9, 1.2);
      \foreach \i in {0,...,5} {
        \pgfmathsetmacro{\ypos}{-0.3 + \i*0.6}
        \node[circle, draw, fill=yellow!20, minimum size=0.35cm, inner sep=0pt] (fc1-\i) at (15.6, \ypos) {};
            }
      \node at (15.6, -1.1) {FC};
      \foreach \i in {0,...,2} {
        \pgfmathsetmacro{\ypos}{0.6 + \i*0.6}
        \node[circle, draw, fill=green!20, minimum size=0.35cm, inner sep=0pt] (fc2-\i) at (17.6, \ypos) {};
            }
      \node at (17.6, -1.1) {Output};
      \foreach \i in {0,...,5} {
        \foreach \j in {0,...,2} {
          \draw[gray!40] (fc1-\i) -- (fc2-\j);
                }
            }
    \end{tikzpicture}
    \caption{Example of a convolutional neural network. An input image passes through convolutional, pooling, and batch normalization layers before being classified by fully connected layers}
    \label{fig:cnn_architecture}
\end{figure}

\begin{definition}[Convolutional Neural Network]
\label{def:cnn}
A convolutional neural network is a neural network $\mathcal{M} = (\mathcal{N}, \mathcal{S})$ with
layers $\mathcal{L} = \{l_1, \dots, l_L\}$, where at least one layer $l \in \mathcal{L}$ is a
convolutional layer.
The remaining layers typically include pooling, batch normalization, and fully connected layers.
\end{definition}

The following paragraphs describe layer types that exist in convolutional neural networks.

\paragraph{Pooling Layers.}
Similar to the convolution operation, pooling slides a window over the feature map.
Instead of a filter, it applies a fixed operation at each position.
This reduces the spatial dimensions of the feature map, lowering the computational cost for
subsequent layers.
We consider two pooling operations: \textit{max pooling} and \textit{average pooling}.
Max pooling takes the values of the region as input and selects the maximum.
Average pooling computes the mean instead.
While max pooling preserves the strongest activation in each region, average pooling provides a
smoother summary by considering all values equally.
Both pooling operations are shown in Figure~\ref{fig:pooling} for a $2 \times 2$ region.

\begin{figure}[ht!]
    \centering
    \begin{tikzpicture}[scale=0.9, every node/.style={font=\footnotesize}]
    % Max pooling
    \draw[gray!60] (0, 0) rectangle ++(1,1);
    \draw[gray!60] (1, 0) rectangle ++(1,1);
    \draw[gray!60] (0,-1) rectangle ++(1,1);
    \draw[gray!60] (1,-1) rectangle ++(1,1);
    \node at (0.5, 0.5) {1}; \node at (1.5, 0.5) {3};
    \node at (0.5,-0.5) {5}; \node at (1.5,-0.5) {2};
    \draw[->, thick, gray] (2.5, 0) -- (3.5, 0);
    \draw[gray!60] (4, -0.5) rectangle ++(1,1);
    \node at (4.5, 0) {5};
    \node[above] at (1, 1.2) {Max pooling};

    % Average pooling
    \begin{scope}[xshift=7cm]
    \draw[gray!60] (0, 0) rectangle ++(1,1);
    \draw[gray!60] (1, 0) rectangle ++(1,1);
    \draw[gray!60] (0,-1) rectangle ++(1,1);
    \draw[gray!60] (1,-1) rectangle ++(1,1);
    \node at (0.5, 0.5) {1}; \node at (1.5, 0.5) {3};
    \node at (0.5,-0.5) {5}; \node at (1.5,-0.5) {2};
    \draw[->, thick, gray] (2.5, 0) -- (3.5, 0);
    \draw[gray!60] (4, -0.5) rectangle ++(1,1);
    \node at (4.5, 0) {2.75};
    \node[above] at (1, 1.2) {Average pooling};
    \end{scope}
    \end{tikzpicture}
    \caption{Max pooling and average pooling for a $2 \times 2$ region. Max pooling selects the maximum, average pooling computes the mean}
    \label{fig:pooling}
\end{figure}



\paragraph{Batch Normalization.}
As a network grows deeper, the activation values passed between layers can shift during training,
making optimization difficult.
Batch normalization addresses this by normalizing these values across a mini-batch of training
samples~\cite{ioffe2015batch}.

For a given value $x$, the batch mean $\mu$ is subtracted and the result is divided by the standard
deviation $\sigma$, centering the values around zero with unit variance:
\begin{equation}
\hat{x} = \frac{x - \mu}{\sigma}
\end{equation}
This stabilizes training and reduces sensitivity to weight initialization.
It also allows higher learning rates, which leads to faster convergence~\cite{ioffe2015batch}.

\paragraph{Flatten.}
The output of the convolutional and pooling stages consists of multiple channels $c$.
Fully connected layers, however, expect a one-dimensional input vector.
A flatten layer combines all values into a single vector so that it can be passed through fully
connected layers to produce the final output predictions.

Now that the building blocks of convolutional neural networks have been introduced, we turn to a
specific CNN architecture that all models in this work are based on: residual networks.

\section{Residual Networks}
\label{sec:residual_networks}

Adding more layers to a neural network does not always improve performance.
In theory, a deeper network should do at least as well as a shallower one, since the extra layers
could just learn the identity function and pass the input through unchanged.
In practice, however, standard networks struggle to learn this identity mapping~\cite{he2016deep}.
Learning the identity requires all weights to reproduce the input exactly, which is difficult to
achieve for every possible input.
He et al.\ proposed residual networks (ResNets) to address this~\cite{he2016deep}.

The key idea is to add skip connections that bypass one or more layers and add the original input
$x$ directly to the output.
Let $H(x)$ denote the desired output of a group of layers. In a ResNet, the layers do not learn
$H(x)$ directly. Instead, they learn the residual function, which is the difference between the
desired output and the input:
\begin{equation}
        F(x) = H(x) - x
\end{equation}
This can be rearranged into the core equation of a residual block, where $F(x)$ is computed by the
main path and $x$ is passed through the skip connection:
\begin{equation}
        H(x) = F(x) + x
\end{equation}
Assuming the optimal mapping is the identity, i.e., $H(x) = x$, this would mean that the result of
the main path $F(x)$ is zero. Pushing all weights towards zero is much easier than learning the
identity mapping $H(x) = x$ from scratch.
The main path and the skip connection together form a residual block. The structure of a single
residual block is shown in Figure~\ref{fig:resnet}.

\begin{figure}[ht!]
    \centering
    \begin{tikzpicture}[
                scale=0.9,
                every node/.style={font=\small},
                block/.style={rectangle, draw, thick, minimum width=2.5cm, minimum height=0.8cm},
                addnode/.style={circle, draw, dashed, thick, minimum size=0.8cm, inner sep=0pt},
        ]
      \node[font=\normalsize] (x) at (0, 0) {$x$};
      \draw[->, thick] (0, -0.3) -- (0, -1.0);
      \node[block] (conv1) at (0, -1.5) {\texttt{Conv}};
      \draw[->, thick] (0, -1.9) -- (0, -2.6);
      \node[block] (relu) at (0, -3.1) {\texttt{ReLU}};
      \draw[->, thick] (0, -3.5) -- (0, -4.2);
      \node[block] (conv2) at (0, -4.7) {\texttt{Conv}};
      \draw[->, dashed, thick] (0, -5.1) -- (0, -6.0);
      \draw[dashed, thick, rounded corners=4pt]
                (-1.8, -0.8) rectangle (1.8, -5.5);
      \node[font=\normalsize] at (-2.4, -3.1) {$F(x)$};
      \node[addnode] (add) at (0, -6.5) {\textbf{+}};
      \node[font=\normalsize, anchor=east] at (-0.7, -6.5) {$y$};
      \draw[->, thick] (0, -6.9) -- (0, -7.5);
      % skip connection
      \draw[thick] (0, -0.3) -- (3.5, -0.3);
      \draw[thick] (3.5, -0.3) -- (3.5, -6.5);
      \draw[->, thick] (3.5, -6.5) -- (0.45, -6.5);
      \node[font=\normalsize] at (5.0, -3.1) {identity ($x$)};
    \end{tikzpicture}
    \caption{A residual block with a skip connection. The input $x$ is added to the output of the block $F(x)$, yielding $y = F(x) + x$}
    \label{fig:resnet}
\end{figure}

The figure shows a simplified residual block with two convolutional layers and a ReLU activation in
between. The left side represents the \textit{main path}, which computes the residual function
$F(x)$. The right side is the \textit{skip connection}, which passes the input $x$ unchanged. Both
paths are combined at the addition node, producing the output $y = F(x) + x$, which corresponds to
our desired mapping $H(x)$.
In practice, each convolutional layer is followed by batch normalization, and a ReLU is applied
after the addition.

In this work, we define a residual block as follows.

\begin{definition}[Residual Block]
\label{def:residual_block}
A residual block $b = \{l_i, \dots, l_j\} \subset \mathcal{L}$ is a group of consecutive layers
that consists of two paths.
The \textit{main path} transforms the input $\mathbf{x}$ through the sequence of layers and
computes the residual function $F(\mathbf{x})$.
The \textit{skip connection} passes the input $\mathbf{x}$ unchanged.
The output of the block is the sum of both paths: $\mathbf{y} = F(\mathbf{x}) + \mathbf{x}$.
We denote the set of all residual blocks in a network as $\mathcal{B}$.
\end{definition}

There are two types of skip connections.
The \textit{identity shortcut} passes the input $x$ through unchanged, as shown in
Figure~\ref{fig:resnet}. This is used when the input and output dimensions are the same.
When the dimensions differ, a \textit{projection shortcut} is used instead. It applies a
transformation to $x$ to match the output dimensions before the addition.

Besides the different types of skip connections, there are also different types of residual
blocks~\cite{he2016deep}.
A \textit{BasicBlock} consists of two convolutional layers, each followed by batch normalization,
with a ReLU activation between them. It is used in shallower residual networks.
A \textit{Bottleneck} block is larger and consists of three convolutional layers, each followed by
batch normalization. The name comes from its structure: the first layer reduces the channel
dimension, the computation is performed on this smaller representation, and the third layer
restores the original dimension. It is used in deeper residual networks.

By stacking multiple residual blocks, a residual network is formed. We define it as follows.

\begin{definition}[Residual Network]
\label{def:resnet}
A residual network is a convolutional neural network $\mathcal{M} = (\mathcal{N}, \mathcal{S})$
whose layers $\mathcal{L}$ are organized into a sequence of residual blocks $\mathcal{B} = \{b_1,
\dots, b_K\}$.
The blocks in $\mathcal{B}$ are either all BasicBlocks or all BottleneckBlocks.
\end{definition}

He et al.~\cite{he2016deep} presented several ResNet variants, named by their total number of
layers.
Table~\ref{tab:resnet_variants} shows the complete list of variants.
Shallower models like ResNet-18 and ResNet-34 use BasicBlocks, while deeper models starting from
ResNet-50 use BottleneckBlocks.
The blocks are organized into four groups, each operating at a different spatial resolution.

\begin{table}[ht]
    \centering
    \begin{tabular}{lccc}
        \toprule
                Model & Block Type & Blocks per Group & Total Layers \\
        \midrule
                ResNet-18 & BasicBlock & 2, 2, 2, 2 & 18 \\
                ResNet-34 & BasicBlock & 3, 4, 6, 3 & 34 \\
                ResNet-50 & Bottleneck & 3, 4, 6, 3 & 50 \\
                ResNet-101 & Bottleneck & 3, 4, 23, 3 & 101 \\
                ResNet-152 & Bottleneck & 3, 8, 36, 3 & 152 \\
        \bottomrule
    \end{tabular}
    \caption{ResNet variants with their block types and configurations~\cite{he2016deep}. Shallower models use BasicBlocks, deeper models use BottleneckBlocks}
    \label{tab:resnet_variants}
\end{table}

With these definitions, we have introduced all network structures and notation relevant for the
following chapters. In this work, we exclusively focus on residual networks.
Besides understanding the network architecture, our approach also builds on the concept of program
slicing, which we introduce in the next section.

\section{Program Slicing}

Program slicing is a technique for extracting relevant parts of a program with respect to a
specific computation~\cite{weiser1981program, weiser1984program}.
Given a variable of interest at a specific point in the program, a slice contains only those
statements that may affect the value of that variable.
We formalize this as follows.

\begin{definition}[Program]
\label{def:program}
A program $P = \{s_1, s_2, \dots, s_n\}$ is a set of statements, where each statement $s_i$
represents an individual instruction.
\end{definition}

Since a program consists of multiple statements, we can define a slice as a subset of these
statements that are relevant to a specific computation.

\begin{definition}[Program Slice]
\label{def:program_slice}
Let $P$ be a program and $C = (s, V)$ a slicing criterion, where $s$ is a statement in $P$ and $V$
is a set of variables.
A program slice $P_C \subseteq P$ is the set of all statements in $P$ that may affect the values of
variables in $V$ at statement $s$.
\end{definition}

Two main variants of program slicing are distinguished in the literature:
\begin{itemize}
    \item \textbf{Static slicing:} Computes slices by analyzing the source code without executing it.
        The resulting slice is valid for all possible inputs to the
    program~\cite{weiser1984program}.
    \item \textbf{Dynamic slicing:} Computes a slice dynamically at runtime during program execution.
        The resulting slice is based on a specific input and only contains statements that were
    actually executed~\cite{korel1988dynamic}.
\end{itemize}
We illustrate both program slicing variants using the example program $P$ in
Listing~\ref{lst:slicing-example}. 
The program consists of eight statements in total. It reads an input value, performs a conditional
computation and finally prints the result.
The slicing criterion is $C = (8, \{\texttt{z}\})$, where $s = 8$ refers to the last statement
\texttt{print(z)} as the point of interest.
The variable we want to observe is $V = \{\texttt{z}\}$.
Both variants then determine which statements in the program affect the value of \texttt{z} at that
point.

\begin{lstlisting}[
        language=Python,
        escapechar=|,
        basicstyle=\ttfamily\small,
        frame=single,
        caption={Example program $P$ with slicing criterion $C = (8, \{\texttt{z}\})$.},
        label={lst:slicing-example}
]
x = input()
y = 5
w = 10
if x > 0:
        z = x + y
else:
        z = x - y
print(z)
\end{lstlisting}

For the static slice, the source code is analyzed without executing it.
Since \texttt{x} is read from user input, its value is unknown and could be any possible value.
Therefore, both the \texttt{if}-branch where $x > 0$ and the \texttt{else}-branch where $x \leq 0$
must be considered.
However, line~3 (\texttt{w~=~10}) does not affect \texttt{z} and is excluded.
The static slice therefore includes all statements in $P$ except for line~3: $P_C =
\{1,2,4,5,6,7,8\}$.

\begin{lstlisting}[
        language=Python,
        escapechar=|,
        basicstyle=\ttfamily\small,
        frame=single,
        caption={Static slice $P_C = \{1,2,4,5,6,7,8\}$ for $C = (8, \{\texttt{z}\})$.},
        label={lst:static-slice}
]
|\colorbox{staticslice}{x = input()}|
|\colorbox{staticslice}{y = 5}|
w = 10
|\colorbox{staticslice}{if x > 0:}|
|\colorbox{staticslice}{~~~~z = x + y}|
|\colorbox{staticslice}{else:}|
|\colorbox{staticslice}{~~~~z = x - y}|
|\colorbox{staticslice}{print(z)}|
\end{lstlisting}

A dynamic slice considers a specific execution with a known input such as $x = 3$.
Since the input is known, only the statements that were actually executed and affected \texttt{z}
are included.
With $x = 3 > 0$, only the \texttt{if}-branch is taken, and both line~3 and the
\texttt{else}-branch are excluded.
The dynamic slice is therefore the following set of statements: $P_C = \{1,2,4,5,8\}$.

\begin{lstlisting}[
        language=Python,
        escapechar=|,
        basicstyle=\ttfamily\small,
        frame=single,
        caption={Dynamic slice $P_C = \{1,2,4,5,8\}$ for $C = (8, \{\texttt{z}\})$ with input $x =
    3$.},
        label={lst:dynamic-slice}
]
|\colorbox{dynamicslice}{x = input()}|
|\colorbox{dynamicslice}{y = 5}|
w = 10
|\colorbox{dynamicslice}{if x > 0:}|
|\colorbox{dynamicslice}{~~~~z = x + y}|
else:
        z = x - y
|\colorbox{dynamicslice}{print(z)}|
\end{lstlisting}

Compared to the static slice, the dynamic slice is smaller because it only considers the execution
path that was actually taken for the given input $x = 3$.
In general, dynamic slices are subsets of static slices and provide a more precise view of the
program behavior for a specific execution.

Over the years, various application areas for program slicing on traditional programs have been
explored.
For example, dynamic slicing can be used in debugging to locate faulty statements by narrowing down
the code that could have caused an incorrect value~\cite{agrawal1993debugging}.
Similarly, slicing has been applied to testing, where it can reduce test suites by identifying and
removing test cases that exercise the same program paths~\cite{arlt2014reducing}.

The concept of dynamic slicing has also been extended beyond traditional software. Zhang et
al.~\cite{zhang2020dynamic} applied it to neural networks.
Instead of a program $P$ with a set of statements, we now consider a neural network $\mathcal{M} =
(\mathcal{N}, \mathcal{S})$ with a set of neurons and synapses.
The slicing criterion also changes: instead of a statement and a set of variables, it consists of a
set of input samples $I$ and a set of output neurons $O$.
The resulting slice is the subset of neurons and synapses that significantly contribute to the
output.
A neural network slice is defined in Definition~\ref{def:nn_slice}. We adopt this definition for
the rest of this work.


\begin{definition}[Neural Network Slice~\cite{zhang2020dynamic}]
\label{def:nn_slice}
Let $\mathcal{M} = (\mathcal{N}, \mathcal{S})$ be a neural network and $C = (I, O)$ a slicing
criterion, where $I = \{\xi_1, \xi_2, \dots, \xi_m\}$ is a set of input samples and $O = \{o_1,
o_2, \dots, o_k\}$ is a set of output neurons of interest.
The neural network slice $\mathcal{M}_C = (\mathcal{N}_C, \mathcal{S}_C)$ with $\mathcal{N}_C
\subseteq \mathcal{N}$ and $\mathcal{S}_C \subseteq \mathcal{S}$ is the subset of neurons and
synapses that significantly contribute to the value of any output $o \in O$ for any input $\xi \in
I$.
\end{definition}

For example, given an input image $\xi$ of a truck and the output neuron $o$ for the class
\textit{truck}, the slice $\mathcal{M}_C$ identifies which neurons and synapses in the network were
responsible for classifying $\xi$ as a truck.

\medskip

In this chapter, we introduced the foundations of neural networks and presented program slicing as
a technique for identifying relevant parts in traditional programs.
We showed how the concept extends to neural networks, where slices identify neurons and synapses
relevant to a specific prediction.
The definitions and notation established here form the basis for the following chapters.
The next chapter reviews related work in program slicing, neural network slicing, and model
interpretability.


\chapter{Related Work}
\label{ch:related-work}

This chapter reviews existing work related to our thesis.
We begin with approaches that apply slicing techniques to neural networks.
After that, we discuss interpretability methods that aim to explain model predictions.
We conclude with program analysis techniques that have been adapted to test and analyze neural
networks.

\section{Neural Network Slicing}
\label{sec:rw-slicing}

Program slicing has recently been adapted to deep neural
networks~\cite{zhang2020dynamic,wang2018interpret,qiu2019adversarial,zhou2024neusemslice}.
The goal shifts from extracting traditional program statements that affect a variable of interest
to identifying the internal components of a network that contribute to a particular output.

NNSlicer applies dynamic program slicing to deep neural networks~\cite{zhang2020dynamic}.
In a profiling phase, it estimates the average activation of each neuron over the training data.
A forward analysis then computes activation deltas for a specific input.
Finally, a backward analysis propagates these deltas from the output layer back through the network
to obtain a slice.
The authors demonstrate that slices can support tasks such as detecting adversarial inputs, pruning
models, and protecting specific predictions.
However, NNSlicer operates exclusively at the individual neuron level.
Modern convolutional neural networks are organized into higher-level structures such as channels
and residual blocks.
Our work extends NNSlicer by aggregating neuron-level information to channels and residual blocks.
We perform slicing at these coarser architectural granularities, which have not been explored in
prior work.

Other work extracts critical paths through a network instead of explicit slices.
Qiu et al.\ propose effective paths, which treat the network as a data-flow
graph~\cite{qiu2019adversarial}.
For each input, they select the minimal set of neurons and synapses whose positive contributions
exceed a fixed fraction of the output logit.
Only positive contributions are considered, meaning that neurons which suppress the output are
excluded.
By aggregating per-input paths into class-level paths, they show that different classes activate
distinct sub-networks and use path similarity for adversarial example detection.
In contrast, our approach builds on NNSlicer, which considers both positive and negative deviations
from the mean activation.
Similar to NNSlicer, effective paths operate exclusively at the neuron level, whereas our method
introduces a hierarchical block and channel structure above the neuron level.

Wang et al.\ interpret neural networks by identifying critical data routing paths through learned
control gates~\cite{wang2018interpret}.
Each gate is a non-negative scalar attached to a channel, where a high value indicates a critical
channel and a low value indicates a less important one.
This is conceptually similar to our channel-level deltas, which also assign a single value per
channel to determine its relevance.
However, the mechanisms differ.
Their gates must be trained separately for each individual input, which adds substantial
computational cost.
Our deltas, in contrast, are derived directly from existing activations without additional
training.

A different line of work targets model maintenance rather than interpretability.
NeuSemSlice introduces semantic slicing, which identifies critical neurons based on contribution
scores~\cite{zhou2024neusemslice}.
For each category and layer, it selects a subset of neurons whose representations retain high
centered-kernel-alignment (CKA) similarity to the full layer.
The selected neurons are then merged across categories into semantic components that support
maintenance tasks such as pruning, fine-tuning, and incremental development.
Unlike our approach, NeuSemSlice measures neuron relevance through CKA similarity across the
training set, whereas we compute per-input activation deltas relative to a profiled mean.
Overall, NeuSemSlice is designed for semantic-aware model maintenance, whereas our work focuses on
efficient slicing of residual networks.



\section{Model Interpretability}
\label{sec:rw-interpretability}


Explainable artificial intelligence (XAI) aims to explain predictions of deep neural networks and
make their decisions more transparent to humans~\cite{Samek_2021,samek2017understanding}.
Nazir et al.\ provide a recent taxonomy of feature attribution techniques and distinguish between
model-agnostic and model-specific approaches~\cite{nazir2025survey}.
Model-agnostic methods treat the network as a black box and operate only on its inputs and outputs.
Model-specific methods, on the other hand, leverage the internal structure and parameters of the
network and can therefore provide explanations that are based on internal computations.
Our work falls into the second category, as we trace activation patterns through the internal
layers of convolutional networks to extract class-specific slices.

Among model-specific methods, reference-based approaches explain predictions by comparing neuron
activations against a reference input.
DeepLIFT computes contribution scores for each neuron by measuring the difference between its
activation and a reference activation~\cite{shrikumar2017deeplift}.
These scores are then propagated backward through the network to quantify how much each neuron
contributed to the output.
This neuron-level perspective is conceptually close to our approach, as we also compute per-neuron
activation differences.
However, DeepLIFT requires a fixed reference input and explains individual predictions, whereas we
aggregate activations over many samples of a class to obtain a class-level view.

Rather than comparing against a reference, activation-based methods directly use the feature maps
produced during inference.
Grad-CAM computes class-specific localization maps by weighting the feature maps of a convolutional
layer with the gradients flowing back from the target class~\cite{selvaraju2017gradcam}.
The resulting heatmap highlights which spatial regions of the input are most relevant for a given
prediction.
Our approach is also activation-based, as we collect neuron activations via forward hooks during
inference.
Instead of producing spatial heatmaps, we aggregate these activations to identify which neurons,
channels, and residual blocks are critical for a given class.

A third line of work propagates relevance scores rather than activations or gradients.
Layer-wise relevance propagation (LRP) redistributes the prediction score backward through the
network by applying local propagation rules at each layer~\cite{bach2015lrp}.
This is conceptually similar to our backward analysis, where we also apply operation-specific rules
layer by layer to propagate the contribution of each neuron to the final slice.
However, LRP decomposes a single prediction down to the input features, while our backward pass
traces contributions across neurons, channels, and blocks to extract class-specific slices.


\section{Program Analysis for Neural Networks}
\label{sec:rw-analysis}

Beyond slicing and interpretability, ideas from traditional program analysis have been applied to
neural networks for testing, fault localization, and model repair.
Unlike traditional program analysis, which operates on source code, the following methods work
directly on the trained network by analyzing its internal activations and parameters.

One line of work adapts fault localization techniques to identify neurons that cause incorrect
predictions.
DeepLocalize monitors internal activations during training and flags neurons whose values deviate
from expected ranges~\cite{wardat2021deeplocalize}.
Hashemifar et al.\ transfer spectrum-based fault localization to neural networks and analyze
execution paths to rank suspicious neurons~\cite{hashemifar2025path}.
Both approaches operate at the neuron level and trace information through network layers, which is
structurally similar to our forward and backward analysis.
However, their goal is to locate faulty neurons that cause mispredictions, whereas we identify
neurons that are relevant for a specific class.

Complementary to fault localization, neural network testing applies coverage criteria to measure
how thoroughly a test suite exercises a model.
DeepXplore introduces neuron coverage and uses differential testing across multiple models to
generate inputs that trigger previously inactive neurons~\cite{pei2017deepxplore}.
DeepGauge extends this with finer-grained metrics at the neuron and layer level, such as
k-multisection coverage and boundary coverage~\cite{ma2018deepgauge}.
These methods share the idea of reasoning about individual neuron activations, but they measure
test adequacy rather than class-specific relevance.

A third line of work applies neuron-level modifications to repair undesirable model behavior once
faults or unfairness have been detected.
RUNNER identifies neurons responsible for unfair predictions and adjusts their weights to reduce
bias~\cite{li2024runner}.
Fayyazi et al.\ use counterexample-guided verification to locate and repair fairness violations in
neural networks~\cite{ijcai2025p42}.
While these methods also operate on individual neurons, they modify the network to correct its
behavior.
Our approach does not alter the model but extracts a subset of it as a standalone slice.


\vspace{0.2cm}

The work in this thesis builds on the data-flow-based approach of NNSlicer and extends it with
coarser slicing granularities.
While interpretability methods like DeepLIFT and Grad-CAM explain which input features are
important, our approach answers a complementary question: which internal components of the network
are responsible for a prediction.
Similar to Grad-CAM, we operate at the feature-map level, but instead of highlighting spatial
regions, we identify which channels and blocks contribute to a decision.
In contrast to testing, fault localization, and repair, we use slicing purely as a descriptive tool
to extract reusable, class-specific subnetworks.
To the best of our knowledge, this is the first approach that extends neuron-level slicing to a
hierarchical block-based method.



\chapter{NNSlicer}
\label{ch:nnslicer}

The following chapter presents a detailed overview of the NNSlicer framework that was introduced by
Zhang et al.~\cite{zhang2020dynamic}.
NNSlicer adapts the concepts of dynamic program slicing to neural networks by calculating the
neurons and synapses that contribute most to a specific output.
Understanding this framework is essential, as our block-based slicing approach builds directly on
its pipeline and contribution rules.

NNSlicer consists of three phases to construct a neural network slice
(Definition~\ref{def:nn_slice}).
In the first \textit{profiling phase}, the mean activation value of each neuron is computed over
the entire training dataset~$\mathcal{I}$.
These means establish a baseline of typical neuron behavior across the dataset.
The second \textit{forward analysis phase} calculates the activation values for a specific input
sample~$\xi \in \mathcal{I}$ and compares them to the mean baseline from profiling.
The resulting activation deltas represent how much each neuron reacts to the input.
In the final \textit{backward analysis phase}, the actual slices are constructed.
Contributions are propagated backward from the target output neurons to the input layer.
The result is a slice $\mathcal{M}_C = (\mathcal{N}_C, \mathcal{S}_C)$, i.e., the subset of neurons
and synapses that significantly contribute to the output.
In the following sections, we take a closer look at each phase with visualized examples.

\section{Profiling}

The profiling phase establishes a baseline for neuron behavior.
Recall that each neuron $n \in \mathcal{N}$ produces an activation value~$y_n$ for a given input.
We write $y_n(\xi)$ for the activation of neuron~$n$ on a specific input sample~$\xi$, and define
its mean activation over the training dataset~$\mathcal{I}$ as:
\begin{equation}
    \bar{y}_n = \frac{1}{|\mathcal{I}|} \sum_{\xi \in \mathcal{I}} y_n(\xi)
\end{equation}
This mean $\bar{y}_n$ captures the typical response of each neuron across all inputs and serves as
a reference point for the forward analysis, where input-specific activations are compared against
it.
Since the mean is computed over the entire dataset~$\mathcal{I}$, profiling requires one forward
pass per sample.
However, profiling only needs to be computed once per model and dataset and can be reused for all
subsequent slicing operations.

Figure~\ref{fig:profiling} illustrates profiling for a simple network.
On the left, the pretrained model contains only the learned weights on its connections.
After profiling, each neuron stores its mean activation value computed over the training data.
For example, the first hidden neuron has a mean activation of $\bar{y}_n = 2$, meaning it produces
a mean activation of~$2$ across all training samples in $\mathcal{I}$.

\begin{figure}[ht!]
    \centering
    \begin{tikzpicture}[scale=0.8, every node/.style={font=\scriptsize},
                neuron/.style={circle, draw, minimum size=0.9cm, inner sep=0pt}]
      \node[neuron] (Li1) at (0, 1.2) {};
      \node[neuron] (Li2) at (0, -1.2) {};
      \node[neuron] (Lh1) at (3, 2.4) {};
      \node[neuron] (Lh2) at (3, 0) {};
      \node[neuron] (Lh3) at (3, -2.4) {};
      \node[neuron] (Lo1) at (6, 1.2) {};
      \node[neuron] (Lo2) at (6, -1.2) {};
      \draw[gray!50, dashed] (Li1) -- node[pos=0.3, fill=white, font=\tiny\bfseries, inner sep=1pt, text=gray!60] {2} (Lh1);
      \draw[gray!50, dashed] (Li1) -- node[pos=0.3, fill=white, font=\tiny\bfseries, inner sep=1pt, text=gray!60] {-3} (Lh2);
      \draw[gray!50, dashed] (Li1) -- node[pos=0.15, fill=white, font=\tiny\bfseries, inner sep=1pt, text=gray!60] {1} (Lh3);
      \draw[gray!50, dashed] (Li2) -- node[pos=0.15, fill=white, font=\tiny\bfseries, inner sep=1pt, text=gray!60] {1} (Lh1);
      \draw[gray!50, dashed] (Li2) -- node[pos=0.3, fill=white, font=\tiny\bfseries, inner sep=1pt, text=gray!60] {2} (Lh2);
      \draw[gray!50, dashed] (Li2) -- node[pos=0.3, fill=white, font=\tiny\bfseries, inner sep=1pt, text=gray!60] {-2} (Lh3);
      \draw[gray!50, dashed] (Lh1) -- node[pos=0.3, fill=white, font=\tiny\bfseries, inner sep=1pt, text=gray!60] {3} (Lo1);
      \draw[gray!50, dashed] (Lh1) -- node[pos=0.3, fill=white, font=\tiny\bfseries, inner sep=1pt, text=gray!60] {-1} (Lo2);
      \draw[gray!50, dashed] (Lh2) -- node[pos=0.3, fill=white, font=\tiny\bfseries, inner sep=1pt, text=gray!60] {-2} (Lo1);
      \draw[gray!50, dashed] (Lh2) -- node[pos=0.3, fill=white, font=\tiny\bfseries, inner sep=1pt, text=gray!60] {3} (Lo2);
      \draw[gray!50, dashed] (Lh3) -- node[pos=0.3, fill=white, font=\tiny\bfseries, inner sep=1pt, text=gray!60] {1} (Lo1);
      \draw[gray!50, dashed] (Lh3) -- node[pos=0.3, fill=white, font=\tiny\bfseries, inner sep=1pt, text=gray!60] {1} (Lo2);
      \node[below, font=\footnotesize] at (3, -3.5) {Pretrained Model};
      \draw[->, very thick, black!60] (7.5, 0) -- node[above, font=\footnotesize] {Profiling} (9.5, 0);
      \begin{scope}[xshift=11cm]
      \node[neuron, fill=blue!10] (Ri1) at (0, 1.2) {3};
      \node[neuron, fill=blue!10] (Ri2) at (0, -1.2) {5};
      \node[neuron, fill=yellow!20] (Rh1) at (3, 2.4) {2};
      \node[neuron, fill=yellow!20] (Rh2) at (3, 0) {4};
      \node[neuron, fill=yellow!20] (Rh3) at (3, -2.4) {1};
      \node[neuron, fill=green!15] (Ro1) at (6, 1.2) {7};
      \node[neuron, fill=green!15] (Ro2) at (6, -1.2) {1};
      \draw[gray!50, dashed] (Ri1) -- node[pos=0.3, fill=white, font=\tiny\bfseries, inner sep=1pt, text=gray!60] {2} (Rh1);
      \draw[gray!50, dashed] (Ri1) -- node[pos=0.3, fill=white, font=\tiny\bfseries, inner sep=1pt, text=gray!60] {-3} (Rh2);
      \draw[gray!50, dashed] (Ri1) -- node[pos=0.15, fill=white, font=\tiny\bfseries, inner sep=1pt, text=gray!60] {1} (Rh3);
      \draw[gray!50, dashed] (Ri2) -- node[pos=0.15, fill=white, font=\tiny\bfseries, inner sep=1pt, text=gray!60] {1} (Rh1);
      \draw[gray!50, dashed] (Ri2) -- node[pos=0.3, fill=white, font=\tiny\bfseries, inner sep=1pt, text=gray!60] {2} (Rh2);
      \draw[gray!50, dashed] (Ri2) -- node[pos=0.3, fill=white, font=\tiny\bfseries, inner sep=1pt, text=gray!60] {-2} (Rh3);
      \draw[gray!50, dashed] (Rh1) -- node[pos=0.3, fill=white, font=\tiny\bfseries, inner sep=1pt, text=gray!60] {3} (Ro1);
      \draw[gray!50, dashed] (Rh1) -- node[pos=0.3, fill=white, font=\tiny\bfseries, inner sep=1pt, text=gray!60] {-1} (Ro2);
      \draw[gray!50, dashed] (Rh2) -- node[pos=0.3, fill=white, font=\tiny\bfseries, inner sep=1pt, text=gray!60] {-2} (Ro1);
      \draw[gray!50, dashed] (Rh2) -- node[pos=0.3, fill=white, font=\tiny\bfseries, inner sep=1pt, text=gray!60] {3} (Ro2);
      \draw[gray!50, dashed] (Rh3) -- node[pos=0.3, fill=white, font=\tiny\bfseries, inner sep=1pt, text=gray!60] {1} (Ro1);
      \draw[gray!50, dashed] (Rh3) -- node[pos=0.3, fill=white, font=\tiny\bfseries, inner sep=1pt, text=gray!60] {1} (Ro2);
      \node[below, font=\footnotesize] at (3, -3.5) {Profiled Model};
      \end{scope}
    \end{tikzpicture}
    \caption{Profiling phase. Mean activations $\bar{y}_n$ are computed for every neuron $n \in \mathcal{N}$ over the training dataset $\mathcal{I}$}
    \label{fig:profiling}
\end{figure}

\section{Forward Analysis}
\label{sec:forward_analysis}

Given the profiled model, the forward analysis phase records activation values for a specific input
sample $\xi \in \mathcal{I}$ and compares them to the baseline from profiling.
For each neuron $n \in \mathcal{N}$, the activation delta is computed as
\begin{equation}
    \Delta_n = y_n(\xi) - \bar{y}_n,
\end{equation}
where $y_n(\xi)$ is the activation of neuron~$n$ for the current input sample~$\xi$ and $\bar{y}_n$
is its mean activation from profiling.
These deltas capture how much the response of a neuron deviates from its typical behavior.
They are used in the backward analysis to weight the contribution of each neuron, so that neurons
with stronger reactions to the input have a larger influence on the resulting slice.
Figure~\ref{fig:forward_analysis} illustrates this process for an input $\xi = [1, 2]$.

\begin{figure}[ht!]
    \centering
    \begin{tikzpicture}[scale=0.8, every node/.style={font=\scriptsize},
                neuron/.style={circle, draw, minimum size=0.9cm, inner sep=0pt}]
      \node[neuron, fill=blue!10] (Li1) at (0, 1.2) {3};
      \node[neuron, fill=blue!10] (Li2) at (0, -1.2) {5};
      \node[neuron, fill=yellow!20] (Lh1) at (3, 2.4) {2};
      \node[neuron, fill=yellow!20] (Lh2) at (3, 0) {4};
      \node[neuron, fill=yellow!20] (Lh3) at (3, -2.4) {1};
      \node[neuron, fill=green!15] (Lo1) at (6, 1.2) {7};
      \node[neuron, fill=green!15] (Lo2) at (6, -1.2) {1};
      \draw[gray!50, dashed] (Li1) -- node[pos=0.3, fill=white, font=\tiny\bfseries, inner sep=1pt, text=gray!60] {2} (Lh1);
      \draw[gray!50, dashed] (Li1) -- node[pos=0.3, fill=white, font=\tiny\bfseries, inner sep=1pt, text=gray!60] {-3} (Lh2);
      \draw[gray!50, dashed] (Li1) -- node[pos=0.15, fill=white, font=\tiny\bfseries, inner sep=1pt, text=gray!60] {1} (Lh3);
      \draw[gray!50, dashed] (Li2) -- node[pos=0.15, fill=white, font=\tiny\bfseries, inner sep=1pt, text=gray!60] {1} (Lh1);
      \draw[gray!50, dashed] (Li2) -- node[pos=0.3, fill=white, font=\tiny\bfseries, inner sep=1pt, text=gray!60] {2} (Lh2);
      \draw[gray!50, dashed] (Li2) -- node[pos=0.3, fill=white, font=\tiny\bfseries, inner sep=1pt, text=gray!60] {-2} (Lh3);
      \draw[gray!50, dashed] (Lh1) -- node[pos=0.3, fill=white, font=\tiny\bfseries, inner sep=1pt, text=gray!60] {3} (Lo1);
      \draw[gray!50, dashed] (Lh1) -- node[pos=0.3, fill=white, font=\tiny\bfseries, inner sep=1pt, text=gray!60] {-1} (Lo2);
      \draw[gray!50, dashed] (Lh2) -- node[pos=0.3, fill=white, font=\tiny\bfseries, inner sep=1pt, text=gray!60] {-2} (Lo1);
      \draw[gray!50, dashed] (Lh2) -- node[pos=0.3, fill=white, font=\tiny\bfseries, inner sep=1pt, text=gray!60] {3} (Lo2);
      \draw[gray!50, dashed] (Lh3) -- node[pos=0.3, fill=white, font=\tiny\bfseries, inner sep=1pt, text=gray!60] {1} (Lo1);
      \draw[gray!50, dashed] (Lh3) -- node[pos=0.3, fill=white, font=\tiny\bfseries, inner sep=1pt, text=gray!60] {1} (Lo2);
      \node[below, font=\footnotesize] at (3, -3.5) {Profiled Model};
      \draw[->, very thick, black!60] (7.5, 0) -- node[above, font=\footnotesize] {Forward} node[below, font=\footnotesize] {Analysis} (9.5, 0);
      \begin{scope}[xshift=11cm]
      \node[neuron, fill=blue!10] (Ri1) at (0, 1.2) {$\Delta{=}\text{-}2$};
      \node[neuron, fill=blue!10] (Ri2) at (0, -1.2) {$\Delta{=}\text{-}3$};
      \node[neuron, fill=yellow!20] (Rh1) at (3, 2.4) {$\Delta{=}2$};
      \node[neuron, fill=yellow!20] (Rh2) at (3, 0) {$\Delta{=}\text{-}3$};
      \node[neuron, fill=yellow!20] (Rh3) at (3, -2.4) {$\Delta{=}\text{-}1$};
      \node[neuron, fill=green!15] (Ro1) at (6, 1.2) {$\Delta{=}3$};
      \node[neuron, fill=green!15] (Ro2) at (6, -1.2) {$\Delta{=}\text{-}2$};
      \draw[gray!50, dashed] (Ri1) -- node[pos=0.3, fill=white, font=\tiny\bfseries, inner sep=1pt, text=gray!60] {2} (Rh1);
      \draw[gray!50, dashed] (Ri1) -- node[pos=0.3, fill=white, font=\tiny\bfseries, inner sep=1pt, text=gray!60] {-3} (Rh2);
      \draw[gray!50, dashed] (Ri1) -- node[pos=0.15, fill=white, font=\tiny\bfseries, inner sep=1pt, text=gray!60] {1} (Rh3);
      \draw[gray!50, dashed] (Ri2) -- node[pos=0.15, fill=white, font=\tiny\bfseries, inner sep=1pt, text=gray!60] {1} (Rh1);
      \draw[gray!50, dashed] (Ri2) -- node[pos=0.3, fill=white, font=\tiny\bfseries, inner sep=1pt, text=gray!60] {2} (Rh2);
      \draw[gray!50, dashed] (Ri2) -- node[pos=0.3, fill=white, font=\tiny\bfseries, inner sep=1pt, text=gray!60] {-2} (Rh3);
      \draw[gray!50, dashed] (Rh1) -- node[pos=0.3, fill=white, font=\tiny\bfseries, inner sep=1pt, text=gray!60] {3} (Ro1);
      \draw[gray!50, dashed] (Rh1) -- node[pos=0.3, fill=white, font=\tiny\bfseries, inner sep=1pt, text=gray!60] {-1} (Ro2);
      \draw[gray!50, dashed] (Rh2) -- node[pos=0.3, fill=white, font=\tiny\bfseries, inner sep=1pt, text=gray!60] {-2} (Ro1);
      \draw[gray!50, dashed] (Rh2) -- node[pos=0.3, fill=white, font=\tiny\bfseries, inner sep=1pt, text=gray!60] {3} (Ro2);
      \draw[gray!50, dashed] (Rh3) -- node[pos=0.3, fill=white, font=\tiny\bfseries, inner sep=1pt, text=gray!60] {1} (Ro1);
      \draw[gray!50, dashed] (Rh3) -- node[pos=0.3, fill=white, font=\tiny\bfseries, inner sep=1pt, text=gray!60] {1} (Ro2);
      \node[below, font=\footnotesize] at (3, -3.5) {Forward Result};
      \end{scope}
    \end{tikzpicture}
    \caption{Forward analysis phase. Activation deltas $\Delta_n = y_n(\xi) - \bar{y}_n$ are computed for a specific input $\xi = [1, 2]$}
    \label{fig:forward_analysis}
\end{figure}

The profiled model on the left shows the mean activations from Figure~\ref{fig:profiling}.
After processing the input, the activation deltas are computed by subtracting the mean values from
the observed activations.
For example, the first hidden neuron has an observed activation of~$4$ and a mean activation
of~$2$, yielding a delta of $\Delta = 4 - 2 = 2$.
A positive delta means the neuron is more active than the mean across the dataset, while a negative
delta means it is less active.
In both cases, a larger absolute value indicates a stronger deviation from the baseline.
The absolute value of the delta $|\Delta_n|$ is called its magnitude and can be interpreted as the
sensitivity of neuron~$n$ to the input.

\section{Backward Analysis}
\label{sec:backward_analysis}

The backward analysis is the final phase, in which the contribution of each neuron and synapse to
the target output neurons~$O$ is computed.
Note that the delta of a neuron is not the same as its contribution to a specific output.
A neuron may react strongly to an input without contributing to the target prediction.
For example, a neuron that detects wing shapes may have a high delta for a given image of a bird.
However, when computing the slice for the output class automobile on this image, the activation of
the neuron does not influence the automobile prediction and has no contribution.

To compute these contributions, NNSlicer considers each neuron $n$ and its $k$ predecessor neurons
$n_1, \dots, n_k$ that are connected to $n$ via synapses.
For each predecessor, a contribution value is computed that describes how much it influences $n$.
NNSlicer distinguishes between two types of contributions:
\begin{itemize}
    \item The \textit{local contribution} $\textit{contrib}_i$ is generated by a single operation and describes how much a preceding neuron~$n_i$ influences the current neuron~$n$ in that computation.
    \item The \textit{cumulative contribution} $\textit{CONTRIB}_n$ accumulates all local contributions that neuron~$n$ receives during the backward pass, representing its overall importance for the target output.
\end{itemize}
Initially, the cumulative contribution of the target output neuron $o \in O$ is set to
$\textit{CONTRIB}_o = 1$.
The backward phase then recursively processes the network from the output layer back to the input
layer.
At each layer $l$, local contributions $\textit{contrib}_i$ are calculated for all predecessor
neurons $n_i \in \mathcal{N}_{l-1}$.
These are accumulated into $\textit{CONTRIB}$, and neurons with $\textit{CONTRIB}_{n_i} \neq 0$
become the new target set.
This process repeats until no more predecessors remain.
Algorithm~\ref{alg:compute_contrib} describes the general workflow.

\begin{algorithm}[ht]
\caption{ComputeContrib~\cite{zhang2020dynamic}: computing the contributions of neurons and synapses for an input sample}\label{alg:compute_contrib}
\begin{algorithmic}[1]
\Require Neural network $\mathcal{M} = (\mathcal{N}, \mathcal{S})$, input sample $\xi$, target neurons $O$
\If{$O$ is empty}
    \State \textbf{terminate}
\EndIf
\State Initialize $O' \gets \emptyset$
\For{each neuron $o \in O$}
    \State Find $o$'s preceding neurons $(\mathcal{N}', \mathcal{S}')$
    \State Compute local contributions of $\mathcal{N}'$ and $\mathcal{S}'$ as $\textit{contrib}$
    \State Update $\textit{CONTRIB}$ with $\textit{contrib}$
\EndFor
\For{each neuron $n \in \mathcal{N}$}
    \State Add $n$ to $O'$ if $n$ is a predecessor of $O$ and $\textit{CONTRIB}_n \neq 0$
\EndFor
\State Obtain $\mathcal{N}'$ by removing neurons in $O$ from $\mathcal{N}$
\State \textbf{call} \textsc{ComputeContrib} with $O \gets O'$ and $\mathcal{N} \gets \mathcal{N}'$
\State \Return $\textit{CONTRIB}$
\end{algorithmic}
\end{algorithm}

\noindent The computation of the local contribution $\textit{contrib}_i$ depends on the type of layer.
Different layers perform different operations, and each requires its own contribution rule.
For the weighted sum operation, which is used in both convolutional and fully connected layers, the
local contribution of the $i$-th predecessor to a neuron~$n$ is:
\begin{equation}
    \textit{contrib}_i = \textit{CONTRIB}_n \cdot \Delta y \cdot w_i \Delta x_i
\end{equation}
where $\Delta y$ is the activation delta of neuron~$n$, $\Delta x_i$ the activation delta of the
preceding neuron~$n_i$, and $w_i$ the synapse weight.
The product $\Delta y \cdot w_i \Delta x_i$ captures the impact that predecessor $n_i$ has on $n$.
For example, if $\Delta y$ is negative and $w_i \Delta x_i$ is positive, then $n_i$ reinforces the
negative deviation of $n$, yielding a contribution that is opposite in sign to
$\textit{CONTRIB}_n$.

Next to the operation rule of weighted sum, NNSlicer defines four more contribution rules for
common neural network operations. 
All supported layers and their corresponding contribution rules are listed in
Table~\ref{tab:operator_rules}.

\begin{table}[ht]
\centering
\caption{Backward contribution rules for each operator type. $\textit{CONTRIB}_n$ is the cumulative contribution of neuron~$n$. $\Delta y$, $\Delta x_i$ are the activation deltas of the output and $i$-th input neuron, respectively~\cite{zhang2020dynamic}}
\label{tab:operator_rules}
\small
\renewcommand{\arraystretch}{1.6}
\begin{tabular}{@{} l l l l @{}}
\toprule
\textbf{Operation} & \textbf{Usage} & \textbf{Definition} & \textbf{Contribution of $i$-th input} \\
\midrule
Weighted sum & Conv./FC layers & $y = \sum_{i=1}^{k} w_i x_i$ & $\textit{CONTRIB}_n \cdot \Delta y
\cdot w_i \Delta x_i$ \\
\hline
Average & Avg-pooling & $y = \frac{1}{k} \sum_{i=1}^{k} x_i$ & $\textit{CONTRIB}_n \cdot \Delta y
\cdot \Delta x_i$ \\
\hline
Maximum & Max-pooling & $y = \max_{i=1}^{k} x_i$ & $\textit{CONTRIB}_n \cdot \Delta y \cdot \Delta
x_i$ if $x_i = y$, else $0$ \\
\hline
Rectify & ReLU & $y = \max(x, 0)$ & $\textit{CONTRIB}_n \cdot \Delta y \cdot \Delta x$ if $x > 0$,
else $0$ \\
\hline
Scale & Batch norm. & $y = \frac{x - \mu}{\sigma}$ & $\textit{CONTRIB}_n \cdot \Delta y \cdot
\Delta x$ \\
\bottomrule
\end{tabular}
\renewcommand{\arraystretch}{1.0}
\end{table}

ReLU, batch normalization, and pooling layers have a simpler structure than convolutional and fully
connected layers.
They each take a single input and apply a fixed transformation.
Their contribution rules reflect this by directly passing the contribution through, scaled by the
activation deltas.

Since different layer types produce contributions at different scales, local contributions can vary
greatly in magnitude.
A fully connected layer with large weights may produce much larger contributions than a
normalization layer.
To prevent any single value from dominating the cumulative result, NNSlicer accumulates only the
sign of each local contribution:
\begin{equation}
    \textit{CONTRIB}_{n_i} \mathrel{+}= \text{sign}(\textit{contrib}_i), \qquad
    \textit{CONTRIB}_{s_i} \mathrel{+}= \text{sign}(\textit{contrib}_i)
\end{equation}
In this way, $\textit{CONTRIB}$ reflects how consistently a neuron contributes across operations,
rather than how strongly it contributes in any single one.
A neuron is considered important if $\textit{CONTRIB}_n \neq 0$, and it may contribute positively
or negatively to the slicing criterion.

Accumulating these contributions across a network with many neurons can be time-consuming.
NNSlicer reduces this cost by excluding neurons with small contributions.
It introduces a threshold parameter~$\theta$ that controls how many local contributions are
included.
Contributions are sorted by magnitude, and the smallest ones are excluded as long as their
cumulative influence on the activation of the neuron remains below~$\theta$.


\begin{figure}[ht!]
    \centering
    \begin{tikzpicture}[scale=1.0, every node/.style={font=\small},
                neuron/.style={circle, draw, minimum size=1.0cm, inner sep=0pt},
                sliced/.style={circle, draw, thick, fill=orange!25, minimum size=1.0cm, inner
        sep=0pt}]
      \node[sliced] (i1) at (0, 1.2) {\tiny $C{=}2$};
      \node[below=0.05cm, font=\tiny, text=gray] at (i1.south) {$\Delta{=}\text{-}2$};
      \node[neuron, fill=gray!10] (i2) at (0, -1.2) {\tiny $C{=}0$};
      \node[below=0.05cm, font=\tiny, text=gray] at (i2.south) {$\Delta{=}\text{-}3$};
      \node[sliced] (h1) at (3.5, 2.4) {\tiny $C{=}1$};
      \node[below=0.05cm, font=\tiny, text=gray] at (h1.south) {$\Delta{=}2$};
      \node[sliced] (h2) at (3.5, 0) {\tiny $C{=}1$};
      \node[below=0.05cm, font=\tiny, text=gray] at (h2.south) {$\Delta{=}\text{-}3$};
      \node[neuron, fill=gray!10] (h3) at (3.5, -2.4) {\tiny $C{=}\text{-}1$};
      \node[below=0.05cm, font=\tiny, text=gray] at (h3.south) {$\Delta{=}\text{-}1$};
      \node[sliced] (o1) at (7, 1.2) {\tiny $C{=}1$};
      \node[below=0.05cm, font=\tiny, text=gray] at (o1.south) {$\Delta{=}3$};
      \draw[->, thick] (8.0, 1.2) -- (7.6, 1.2);
      \node[font=\footnotesize, right] at (8.0, 1.2) {target};
      \node[neuron, fill=gray!10] (o2) at (7, -1.2) {\tiny $C{=}0$};
      \node[below=0.05cm, font=\tiny, text=gray] at (o2.south) {$\Delta{=}\text{-}2$};
      \draw[orange!80!red, thick, ->] (h1) -- node[pos=0.3, fill=white, font=\tiny\bfseries, inner sep=1pt, text=orange!80!red] {3} (o1);
      \draw[orange!80!red, thick, ->] (h2) -- node[pos=0.3, fill=white, font=\tiny\bfseries, inner sep=1pt, text=orange!80!red] {-2} (o1);
      \draw[orange!80!red, thick, ->] (i1) -- node[pos=0.3, fill=white, font=\tiny\bfseries, inner sep=1pt, text=orange!80!red] {2} (h1);
      \draw[orange!80!red, thick, ->] (i1) -- node[pos=0.3, fill=white, font=\tiny\bfseries, inner sep=1pt, text=orange!80!red] {-3} (h2);
      \draw[gray!30, dashed] (i1) -- node[pos=0.15, fill=white, font=\tiny\bfseries, inner sep=1pt, text=gray!60] {1} (h3);
      \draw[gray!30, dashed] (i2) -- node[pos=0.15, fill=white, font=\tiny\bfseries, inner sep=1pt, text=gray!60] {1} (h1);
      \draw[gray!30, dashed] (i2) -- node[pos=0.3, fill=white, font=\tiny\bfseries, inner sep=1pt, text=gray!60] {2} (h2);
      \draw[gray!30, dashed] (i2) -- node[pos=0.3, fill=white, font=\tiny\bfseries, inner sep=1pt, text=gray!60] {-2} (h3);
      \draw[gray!30, dashed] (h1) -- node[pos=0.3, fill=white, font=\tiny\bfseries, inner sep=1pt, text=gray!60] {-1} (o2);
      \draw[gray!30, dashed] (h2) -- node[pos=0.3, fill=white, font=\tiny\bfseries, inner sep=1pt, text=gray!60] {3} (o2);
      \draw[gray!30, dashed] (h3) -- node[pos=0.3, fill=white, font=\tiny\bfseries, inner sep=1pt, text=gray!60] {1} (o1);
      \draw[gray!30, dashed] (h3) -- node[pos=0.3, fill=white, font=\tiny\bfseries, inner sep=1pt, text=gray!60] {1} (o2);
    \end{tikzpicture}
    \caption{Backward analysis phase. Contributions are traced backward from the target neuron. Orange nodes and edges form the computed slice, gray elements are excluded}
    \label{fig:backward_analysis}
\end{figure}

To illustrate these concepts, we apply the backward analysis to the same example network from the
previous phases. We use a threshold of $\theta = 0.2$ throughout the example.
The resulting slice for the target output neuron~$o_1$ is shown in
Figure~\ref{fig:backward_analysis}, where orange neurons and synapses are included in the slice and
gray ones are excluded.

Following Algorithm~\ref{alg:compute_contrib}, the backward analysis starts by setting
$\textit{CONTRIB}_{o_1} = 1$.
The target output neuron $o_1$ has $\Delta y = 3$ and is connected to three predecessor neurons in
the hidden layer via synapses with weights $w_1 = 3$, $w_2 = -2$, and $w_3 = 1$.
Since the network consists of fully connected layers, the weighted sum contribution rule is
applied.
Before computing the local contributions, the threshold~$\theta$ is applied.
The connection between $o_1$ and the first hidden neuron has a magnitude of $|w_1 \cdot \Delta x_1|
= |3 \cdot 2| = 6$.
For the second hidden neuron, the magnitude is $|w_2 \cdot \Delta x_2| = |(-2) \cdot (-3)| = 6$.
The third has a much smaller magnitude of $|w_3 \cdot \Delta x_3| = |1 \cdot (-1)| = 1$.
The total magnitude across all three connections is $6 + 6 + 1 = 13$.
With $\theta = 0.2$, the smallest contributions can be excluded as long as their share of the total
remains below~$\theta$.
Since the third connection accounts for only $1/13 < 0.2$, it is excluded.
For the two remaining predecessor neurons, the local contributions are computed with the weighted
sum contribution rule.
For the first hidden neuron, the local contribution is $\textit{contrib}_1 = 1 \cdot 3 \cdot 3
\cdot 2 = 18$.
For the second hidden neuron, the local contribution is $\textit{contrib}_2 = 1 \cdot 3 \cdot (-2)
\cdot (-3) = 18$.
The cumulative contributions are then updated by adding the sign of the local contributions:
$\textit{CONTRIB}_{h_1} \mathrel{+}= \text{sign}(18) = +1$ and $\textit{CONTRIB}_{h_2} \mathrel{+}=
\text{sign}(18) = +1$.
Both hidden neurons are added to the slice, while the third hidden neuron with $\textit{CONTRIB} =
0$ is excluded.
The target set is then updated to $O' = \{h_1, h_2\}$ and the process repeats for the next layer.
This continues until all layers have been processed.
The neurons with $\textit{CONTRIB} \neq 0$ and the synapses connecting them form the resulting
slice $\mathcal{M}_C = (\mathcal{N}_C, \mathcal{S}_C)$.

\medskip

This chapter presented the three phases of NNSlicer and illustrated them with a concrete example to
create a slice for a target output neuron.
While NNSlicer operates at the level of individual neurons and synapses, the next chapter presents
our block-based slicing approach and concludes with the vectorized implementation.


\chapter{Block-Based Dynamic Slicing}
\label{ch:approach}

In this chapter, we present our block-based slicing approach, which extends the neuron-level
analysis of NNSlicer to operate at a coarser granularity.
\textit{Block slicing} determines which residual blocks to retain from all blocks in the network. \textit{Channel slicing} then selects which channels to retain within the remaining blocks. Only after both filtering steps, the neuron-level backward analysis is performed on the remaining components.
We extend the NNSlicer workflow with the following steps. They are executed in the listed order,
forming a hierarchical pipeline that works from the coarsest level down to individual neurons:
\begin{enumerate}
    \item \textbf{Forward analysis with delta aggregation:}
        During the forward analysis, activation deltas are recorded for each neuron.
        We aggregate these deltas into channel-level, layer-level, and block-level deltas.
    \item \textbf{Block slicing:}
        The block deltas are used to identify residual blocks that show little response to the
    current input.
        These blocks are skipped during the backward analysis.
    \item \textbf{Channel slicing:}
        Within the retained blocks, channel deltas are used to select the most responsive channels
    in each convolutional layer.
    \item \textbf{Vectorized backward analysis:}
        For the remaining neurons in the retained channels and blocks, the backward analysis is
    performed using vectorized tensor operations.
\end{enumerate}

The chapter is structured as follows. The first section describes the delta aggregation procedure,
which provides the basis for the filtering steps. The next two sections present block slicing and
channel slicing, respectively. The chapter concludes with the vectorized backward analysis.

\section{Delta Aggregation}
\label{sec:delta_aggregation}

The backward analysis is the most computationally expensive phase of NNSlicer, as it recursively
propagates contributions
through neurons layer by layer.
For every individual input sample $\xi$ and output neuron $o$, the backward analysis has to be
repeated across the entire network to create the slice for that target.
As a network grows in size and complexity, the number of neurons increases and the backward
analysis becomes increasingly expensive. 
Our goal is to identify irrelevant components before the backward analysis begins, so that the
computation for those
parts of the network can be skipped entirely.

NNSlicer performs two phases before backward analysis: profiling and forward analysis. 
Profiling computes the mean activation $\bar{y}_n$ for each neuron over the entire
dataset~$\mathcal{I}$.
The computational cost of profiling is also high, as it requires a forward pass for each training
sample $\xi \in \mathcal{I}$ to capture the activations.
This cost increases with the size of the dataset to be profiled.
In contrast to the backward analysis, profiling only needs to be computed once and can be reused
for all inputs.
However, profiling is not suitable as a filtering criterion in our case.
The resulting averages are dataset-wide and do not provide information specific to an individual
input. Since we perform dynamic slicing, we need input-specific values to determine which
components are relevant for a given prediction.

The forward analysis provides exactly this. It computes activation deltas $\Delta_n$ for every
neuron, 
which measure how much the response of the neuron for the current input deviates from its dataset
average.
A large deviation from the average, whether positive or negative, suggests that the neuron reacts
strongly to this particular input.
A delta close to zero means it behaves similar to its average and is unlikely to carry
input-specific information.
We extend this idea to coarser structures.
Since neural networks are organized into channels, layers, and blocks, we can aggregate
neuron-level deltas into a single value per structure.
This gives us a criterion that is available right after the forward pass and before we start the
backward analysis.

The hierarchical delta aggregation is shown in Figure~\ref{fig:delta_aggregation}.
Neuron deltas are aggregated into channel deltas.
Channel deltas are aggregated into layer deltas.
Finally, layer deltas are aggregated into block deltas.


\begin{figure}[ht!]
    \centering
    \begin{tikzpicture}[every node/.style={font=\footnotesize}]
      \draw[thick, rounded corners, fill=green!8] (0, 0) rectangle (9, 5.8);
      \node[font=\footnotesize] at (4.5, 5.5) {$\Delta_{\text{block}}^{(b)}$};
      \draw[thick, rounded corners, fill=red!8] (0.3, 0.3) rectangle (8.7, 5.1);
      \node[font=\footnotesize] at (4.5, 4.75) {$\Delta_{\text{layer}}^{(l)}$};
      \draw[thick, rounded corners, fill=blue!15] (2.9, 0.6) rectangle (7.1, 4.3);
      \draw[thick, rounded corners, fill=blue!12] (2.4, 0.6) rectangle (6.6, 4.3);
      \draw[thick, rounded corners, fill=blue!8] (1.9, 0.6) rectangle (6.1, 4.3);
      \node[font=\footnotesize] at (4.0, 3.95) {$\Delta_{\text{ch}}^{(c)}$};
      \draw[thick, rounded corners, fill=yellow!20] (2.2, 2.5) rectangle (3.9, 3.5);
      \node at (3.05, 3.0) {$\Delta_{n_1}$};
      \draw[thick, rounded corners, fill=yellow!20] (4.2, 2.5) rectangle (5.9, 3.5);
      \node at (5.05, 3.0) {$\Delta_{n_2}$};
      \draw[thick, rounded corners, fill=yellow!20] (2.2, 0.9) rectangle (3.9, 1.9);
      \node at (3.05, 1.4) {$\Delta_{n_3}$};
      \draw[thick, rounded corners, fill=yellow!20] (4.2, 0.9) rectangle (5.9, 1.9);
      \node at (5.05, 1.4) {$\Delta_{n_4}$};
    \end{tikzpicture}
    \caption{Hierarchical delta aggregation. Neuron deltas are aggregated into channel deltas, then into layer deltas, and finally into block deltas}
    \label{fig:delta_aggregation}
\end{figure}

We now define each delta level from bottom to top, starting with channels and ending with blocks.
Recall that a channel $c$ consists of a spatial grid of neurons $\mathcal{N}_{\text{ch}}$ that
share the same kernel weights.
By aggregating the neuron-level deltas within a channel, we obtain a single value that describes
how much the activation of the channel deviates from its mean behavior.
We use the absolute value of the delta because we are interested in how much the behavior of a
neuron deviates from its average, regardless of direction.
A positive delta means the neuron is more active than usual for the given input, while a negative
delta means it is less active.
Both directions indicate that the neuron reacts to the input.
Without the absolute value, positive and negative deviations within the same channel could cancel
each other out, making an active channel appear inactive.

We take the mean rather than the sum to make the channel delta independent of the spatial
dimensions.
This ensures that channels of different sizes can be compared fairly.
For example, in deeper layers of a convolutional network, the spatial dimensions shrink due to
pooling.
Using the sum, channels in earlier layers with larger spatial dimensions would systematically
receive higher values than channels in deeper layers, regardless of how strongly their neurons
react.
We therefore define the channel delta as the mean of the absolute neuron deltas within the channel.

\begin{definition}[Channel Delta]
For a channel $c$ with neurons $\mathcal{N}_{\text{ch}} = \{n_{i,j} \mid 1 \leq i \leq H,\; 1 \leq
j \leq W\}$ arranged in a spatial grid of height $H$ and width $W$, the \textit{channel delta} is
the mean absolute neuron delta over all neurons in the channel:
\begin{equation}
    \Delta_{\text{ch}} = \frac{1}{|\mathcal{N}_{\text{ch}}|} \sum_{n \in \mathcal{N}_{\text{ch}}} |\Delta_{n}|
\end{equation}
\end{definition}

Consider a channel with four neurons whose deltas are $\Delta_n = 0.9, -0.7, 0.8, -0.6$.
The channel delta is $\Delta_{\text{ch}} = \frac{1}{4}(|0.9| + |-0.7| + |0.8| + |-0.6|) = 0.75$.
Note that both positive and negative deviations contribute equally, since the absolute value
ensures they do not cancel out.

The same idea of aggregation can be applied at the layer level.
A convolutional layer $l$ consists of a set of channels $\mathcal{C}$.
The layer delta is the mean over all channel deltas $\Delta_{\text{ch}}^{(c)}$ of its channels $c
\in \mathcal{C}$.
As with the channel delta, we use the mean to ensure that layers with different numbers of channels
are comparable.
For fully connected layers without channels, the layer delta is instead computed directly as the
mean over the absolute neuron deltas $|\Delta_n|$ of its neurons $n \in \mathcal{N}_l$.

\begin{definition}[Layer Delta]
For a convolutional layer $l$ with channels $\mathcal{C} = \{c_1, \dots, c_n\}$, the \textit{layer
delta} is the mean of the channel deltas of its channels:
\begin{equation}
    \Delta_{\text{layer}} = \frac{1}{|\mathcal{C}|} \sum_{c \in \mathcal{C}} \Delta_{\text{ch}}^{(c)}
\end{equation}
For a fully connected layer $l$ with neurons $\mathcal{N}_l = \{n_1, \dots, n_m\}$, the layer delta
is defined as:
\begin{equation}
    \Delta_{\text{layer}} = \frac{1}{|\mathcal{N}_l|} \sum_{n \in \mathcal{N}_l} |\Delta_n|
\end{equation}
\end{definition}
A high layer delta means that many of its channels deviate from their average behavior.
A possible interpretation is that an early layer detecting edges and textures may have a high delta
for images with strong visual patterns but a low delta for uniform backgrounds.

Note that we do not use layer deltas directly as a filtering criterion.
Removing an entire layer from the network would break the data flow, as subsequent layers depend on
its output.
Residual blocks, by contrast, include a skip connection that passes the input unchanged.
This allows skipping a block without interrupting the flow through the network.
Layer deltas serve instead as an intermediate step to aggregate channel deltas into block deltas.

A residual network organizes its layers into a set of residual blocks $\mathcal{B} = \{b_1, \dots,
b_K\}$, where each block $b \in \mathcal{B}$ groups multiple consecutive layers.
We average their layer deltas $\Delta_{\text{layer}}^{(l)}$ for each layer $l \in b$ into a single
block delta.

\begin{definition}[Block Delta]
For a residual block $b = \{l_1, \dots, l_n\}$, the \textit{block delta} is the mean of the layer
deltas of its layers:
\begin{equation}
    \Delta_{\text{block}} = \frac{1}{|b|} \sum_{l \in b} \Delta_{\text{layer}}^{(l)}
\end{equation}
\end{definition}
Each residual block acts as a functional unit within the network.
A large block delta indicates that this unit as a whole is active for the current input.
For instance, a block that is sensitive to object shapes could show a high delta for object images
but remain close to zero for scenes dominated by textures.

With these aggregated deltas, we now have an indicator for both channels and blocks that suggests
which components may be relevant for a given input.
The remaining question is how to decide which components to retain before the backward analysis
starts.
The following two sections describe this selection process in detail: first at the block level,
then at the channel level.

\section{Block Slicing}
\label{sec:block_slicing}

The idea behind block slicing is to slice the network at its coarsest granularity before the
backward analysis begins.
As described in the previous section, blocks with high deltas indicate strong reactions to the
current input, while blocks with deltas close to zero are candidates for skipping.
To select which blocks to retain, we adopt a cumulative energy filtering strategy.
We define the total energy as the sum of all block deltas.

\begin{definition}[Block Energy]
The \textit{block energy} is the sum of block deltas over all blocks $\mathcal{B}$:
\begin{equation}
        E_{\text{block}} = \sum_{b \in \mathcal{B}} \Delta_{\text{block}}^{(b)}
\end{equation}
\end{definition}

This gives us a measure of the overall response of the network to the current input at the block
level.
A block with a high delta contributes a large share to the total energy, while a block with a low
delta contributes less.

We introduce a new hyperparameter $\beta \in [0, 1]$ to control this filtering process.
First, we sort all blocks by their delta in descending order.
We then accumulate deltas starting from the block with the highest delta until the cumulative sum
reaches a fraction~$\beta$ of $E_{\text{block}}$.
All remaining blocks are candidates for skipping.
A lower value of $\beta$ means that less cumulative energy is required, so fewer blocks are
retained and more are skipped.
This implies a monotonicity property: lowering $\beta$ always produces a subset of the blocks
retained at a higher $\beta$, since the algorithm accumulates blocks in the same sorted order.
Setting $\beta$ to $1.0$ requires the full energy to be accumulated, which means all blocks are
retained.
This is equivalent to the original NNSlicer without block filtering.
The strategy itself is inspired by the threshold~$\theta$ in the backward analysis of NNSlicer,
where contributions are similarly sorted by magnitude and retained until a fraction of the total is
reached.
We chose this approach to maintain a consistent filtering strategy across all levels of the
hierarchy.
Algorithm~\ref{alg:block_slicing} formalizes the block slicing process based on cumulative delta
energy.

\begin{algorithm}[ht]
\caption{BlockSlice: identifying residual blocks to skip}\label{alg:block_slicing}
\begin{algorithmic}[1]
\Require Block deltas $\Delta_{\text{block}}^{(b)}$ for each $b \in \mathcal{B}$, total energy $E_{\text{block}}$, threshold $\beta \in [0, 1]$
\Ensure Set of blocks to skip $\textit{skip}$
\State Sort blocks by $\Delta_{\text{block}}^{(b)}$ in descending order
\State $\textit{accumulated} \gets 0$
\State $\textit{keep} \gets \emptyset$
\For{each block $b$ in sorted order}
    \State $\textit{accumulated} \gets \textit{accumulated} + \Delta_{\text{block}}^{(b)}$
    \State Add $b$ to $\textit{keep}$
    \If{$\textit{accumulated} \geq \beta \cdot E_{\text{block}}$}
        \State \textbf{break}
    \EndIf
\EndFor
\State $\textit{skip} \gets \{b \notin \textit{keep} \text{ and } b \text{ has identity shortcut}\}$
\State \Return $\textit{skip}$
\end{algorithmic}
\end{algorithm}

In line~11 of the algorithm, the set of skipped blocks is constructed.
A block is skipped only if it is not in the keep set and has an identity shortcut.
Since the identity shortcut passes the input through unchanged, skipping the main path does not
disrupt the data flow.
Blocks with projection shortcuts are always retained.
In principle, blocks with projection shortcuts could also be skipped.
However, this would require special handling of the dimension changes introduced by the shortcut.
Since there are typically only a few projection shortcuts per network, we focus on identity
shortcuts and leave the extension to projection shortcuts as future work.

Figure~\ref{fig:block_slicing} illustrates block slicing with a concrete example.
Consider a network with four residual blocks whose deltas are $\Delta_{\text{block}}^{(b_1)} =
2.2$, $\Delta_{\text{block}}^{(b_2)} = 3.6$, $\Delta_{\text{block}}^{(b_3)} = 0.4$, and
$\Delta_{\text{block}}^{(b_4)} = 1.8$.
The total energy is the sum of all block deltas, yielding $E_{\text{block}} = 8.0$.
We first sort the blocks by delta in descending order, resulting in the sequence $3.6$, $2.2$,
$1.8$, and $0.4$ for blocks $b_2$, $b_1$, $b_4$, and $b_3$, respectively.
With a threshold of $\beta = 0.8$, the energy target is $\beta \cdot E_{\text{block}} = 6.4$.
The algorithm accumulates delta starting from the highest: after including the first three blocks
($3.6 + 2.2 + 1.8 = 7.6 \geq 6.4$), the cumulative delta exceeds the target.
The remaining block with delta $0.4$ is not retained and, assuming it has an identity shortcut, can
be skipped.
For this skipped block, the backward analysis bypasses all layers $l_1, \dots, l_n$ in the main
path of $b_3$ and propagates contributions solely through the identity shortcut.
The backward contribution flow is shown in orange in Figure~\ref{fig:block_slicing}.

\begin{figure}[ht!]
    \centering
    \begin{tikzpicture}[
                block/.style={draw, rounded corners, minimum width=2.2cm, minimum height=1.4cm,
        align=center, font=\small},
                kept/.style={block, fill=orange!25, draw=black, thick},
                skipped/.style={block, fill=gray!10, draw=black, dashed},
                arrow/.style={->, thick, >=stealth},
                delta/.style={font=\footnotesize, text=gray},
        ]
    \node[kept] (b1) at (0, 0) {$\Delta_{b_1} = 2.2$};

    \node[kept] (b2) at (3.5, 0) {$\Delta_{b_2} = 3.6$};

    \node[skipped] (b3) at (7, 0) {$\Delta_{b_3} = 0.4$};

    \node[kept] (b4) at (10.5, 0) {$\Delta_{b_4} = 1.8$};

    % Connections between blocks (backward flow: right to left)
    \draw[arrow, orange!80!red, thick] (b2.west) -- (b1.east);
    \draw[arrow, gray!50, dashed] (b3.west) -- (b2.east);
    \draw[arrow, gray!50, dashed] (b4.west) -- (b3.east);

    % Identity shortcut over skipped block (from b4 to b2, skipping b3)
    \draw[arrow, orange!80!red, thick] (b4.north west) to[bend right=30] node[above, font=\footnotesize, text=black] {identity shortcut} (b2.north east);
    \end{tikzpicture}
    \caption{Block slicing example with $\beta = 0.8$. Blocks $b_1$, $b_2$, and $b_4$ are retained. Block $b_3$ is skipped through its identity shortcut. Orange elements are part of the slice, gray elements are not.}
    \label{fig:block_slicing}
\end{figure}

Now that we know which blocks to skip, the next question is how contributions flow through the
identity shortcut of a skipped block.
Each residual block ends with an addition $y = F(x) + x$, where $F(x)$ is the main path output and
$x$ is the shortcut input.
During the backward analysis, contributions can flow through both paths.
In the main path $F(x)$, the standard contribution rules are applied layer by layer as defined by
NNSlicer.
For the shortcut, however, no existing rule is yet defined.
As the shortcut passes the input through unchanged without modifying or combining values, we treat
it as an unweighted operation.
This matches the behavior of the passthrough rule for unweighted operations:
\begin{equation}
    \textit{contrib} = \textit{CONTRIB}_n \cdot \Delta y \cdot \Delta x
\end{equation}
where $\Delta y$ is the activation delta of the block from which the shortcut starts and $\Delta x$
is the activation delta of the block to which it leads.
Consider the example above, where $b_3$ is skipped.
The contribution arrives from $b_4$ with cumulative contribution $\textit{CONTRIB}_n$ and
activation delta $\Delta y$.
Since the main path of $b_3$ is bypassed, we propagate this contribution directly to $b_2$ through
the identity shortcut.
With the activation delta $\Delta x$ at the output of $b_2$, the contribution passed to $b_2$ is
$\textit{CONTRIB}_n \cdot \Delta y \cdot \Delta x$.
If $b_3$ were retained, the backward analysis would process all layers in its main path as well.
Both the main path and the shortcut contributions would eventually reach $b_2$, where they are
accumulated.

Propagating contributions through the identity shortcut assumes that we know the block structure of
the network. Where do blocks start and end? Which layers belong the to mainpath of a block?
Modern deep learning frameworks such as PyTorch and TensorFlow define models as compositions of
individual layers~\cite{paszke2019pytorch, abadi2016tensorflow}.
While blocks exist as abstractions, there is no standardized way to query the block structure.
The following paragraph describes our approach to automatically detecting blocks and their shortcut
types.

\paragraph{Block Detection.}
In our reimplementation of NNSlicer, we use the FX tracing module of PyTorch to automatically
detect the block structure~\cite{paszke2019pytorch}.
FX produces a computational graph where each node represents an operation.
A node can be a layer, a function, or a simple arithmetic operation.
Residual blocks follow a specific pattern in this graph: each block ends with a single addition
node where the main path and the shortcut come together.
Since in ResNets addition nodes are exclusively used for the skip connection, our first step is to
locate all addition nodes in the graph.
Each addition node has two parent nodes, one from the main path and one from the shortcut.
The main path parent is the last batch normalization layer of the current block.
The shortcut parent is either the output of the previous block, indicating an identity shortcut, or
a convolutional layer that adjusts dimensions, indicating a projection shortcut.
Based on this distinction, we determine whether a block can be skipped: blocks with identity
shortcuts are candidates for skipping, while blocks with projection shortcuts are always retained.
All nodes between the output of the previous block and the addition node form the main path of the
detected block.

The computational graph for a single residual block is illustrated in
Figure~\ref{fig:block_detection}.
\begin{figure}[ht!]
    \centering
    \begin{tikzpicture}[
                nd/.style={draw, circle, minimum size=1cm, font=\scriptsize, inner sep=1pt},
                mainnd/.style={nd, fill=blue!15},
                addnd/.style={nd, fill=blue!30, thick},
                othernd/.style={nd, fill=gray!10},
                arr/.style={->, thick, >=stealth},
        ]
    % Previous block output
    \node[othernd] (prev) at (0, 0) {ReLU};

    % Main path
    \node[mainnd] (conv) at (2.5, 0) {Conv};
    \node[mainnd] (bn) at (5, 0) {BN};

    % Add node
    \node[addnd] (add) at (7.5, 0) {Add};

    % Next
    \node[othernd] (next) at (10, 0) {ReLU};

    % Main path arrows
    \draw[arr] (prev) -- (conv);
    \draw[arr] (conv) -- (bn);
    \draw[arr] (bn) -- (add);
    \draw[arr] (add) -- (next);

    % Identity shortcut
    \draw[arr, dashed] (prev.south) to[out=-30, in=-150] node[below, font=\scriptsize] {identity shortcut} (add.south);

    % Main path boundary (skippable)
    \draw[dashed, red!60, thick, rounded corners] (1.5, -0.7) rectangle (6.0, 0.7);
    \node[font=\scriptsize, red!60] at (3.75, 0.9) {main path (skip)};

    \end{tikzpicture}
    \caption{FX graph of a residual block. The addition node marks the end of the block. The main path and shortcut type are identified by inspecting the parent nodes of the addition.}
    \label{fig:block_detection}
\end{figure}

The example shows a simplified block with one convolutional layer.
Larger blocks follow the same pattern but contain more convolutional layers in the main path.
In both cases, the addition node marks the end of the block.
Its parent nodes reveal the main path and the shortcut.
The batch normalization (BN) node before the addition marks the end of the main path. The shortcut
side connects directly from the ReLU after the addition node of the previous block, confirming an
identity shortcut.
In case of a skip, the main path nodes (Conv, BN) would be excluded from the backward analysis. The
contribution would instead flow through the identity shortcut from the previous ReLU directly to
the Add node.

In this section, we presented block slicing as a method to filter entire residual blocks before the
backward analysis.
We described how blocks are selected based on cumulative energy thresholding and how the block
structure is detected automatically from the computational graph. Skipped blocks pass contributions
through their identity shortcut.
The next section introduces channel slicing, which operates at a finer granularity within the
retained blocks.

\section{Channel Slicing}
\label{sec:channel_slicing}

While block slicing operates at the coarsest structural level and can skip entire residual blocks,
channel slicing operates at the level of individual channels within convolutional layers.
Each residual block contains convolutional layers that consist of multiple channels.
Our motivation for channel slicing is to not consider all channels during the backward analysis,
but only those that show a strong response to the current input.
The channel deltas $\Delta_{\text{ch}}$ give us an indicator of which channels respond strongly and
which do not.
Similar to block slicing, we use these channel deltas to filter out channels before the backward
analysis begins.
We apply the same cumulative energy thresholding strategy and define the total channel energy as
the sum of all channel deltas $\Delta_{\text{ch}}^{(c)}$ in a layer $l$.

\begin{definition}[Channel Energy]
The \textit{channel energy} for a layer $l$ with channels $\mathcal{C}_l$ is the sum of channel
deltas:
\begin{equation}
        E_{\text{ch}} = \sum_{c \in \mathcal{C}_l} \Delta_{\text{ch}}^{(c)}
\end{equation}
\end{definition}

We introduce a new hyperparameter $\alpha \in [0, 1]$ to control channel filtering.
Channels are sorted by their delta in descending order and retained until their cumulative energy
reaches a fraction~$\alpha$ of $E_{\text{ch}}$.
The procedure is described in Algorithm~\ref{alg:channel_slicing}.

\begin{algorithm}[ht]
\caption{ChannelSlice: identifying channels to skip}\label{alg:channel_slicing}
\begin{algorithmic}[1]
\Require Channel deltas $\Delta_{\text{ch}}^{(1)}, \dots, \Delta_{\text{ch}}^{(C)}$, threshold $\alpha \in [0, 1]$
\Ensure Set of active channels $\textit{active}$
\State Sort channels by $|\Delta_{\text{ch}}^{(k)}|$ in descending order
\State $\textit{total} \gets \sum_{k=1}^{C} |\Delta_{\text{ch}}^{(k)}|$
\State $\textit{accumulated} \gets 0$
\State $\textit{active} \gets \emptyset$
\For{each channel $k$ in sorted order}
    \State $\textit{accumulated} \gets \textit{accumulated} + |\Delta_{\text{ch}}^{(k)}|$
    \State Add $k$ to $\textit{active}$
    \If{$\textit{accumulated} \geq \alpha \cdot \textit{total}$}
        \State \textbf{break}
    \EndIf
\EndFor
\State \Return $\textit{active}$
\end{algorithmic}
\end{algorithm}

As with block slicing, this has a monotonicity property: lowering $\alpha$ always produces a subset
of the channels retained at a higher $\alpha$.
Setting $\alpha = 1.0$ retains all channels and is equivalent to the original NNSlicer without
channel filtering.
The two hyperparameters $\alpha$ and $\beta$ operate independently at different levels of the
network hierarchy.
Setting only $\beta < 1$ filters at the block level while keeping all channels.
Setting only $\alpha < 1$ filters at the channel level while keeping all blocks.
Combining both applies filtering at both levels.
Setting $\alpha = \beta = 1$ retains all blocks and all channels, which produces the same result as
the original NNSlicer.

Channel slicing is applied to every convolutional layer in the retained blocks.
Fully connected layers, which do not have a channel structure, are not affected by channel slicing
and remain subject to the full neuron-level backward analysis.

A concrete example is shown in Figure~\ref{fig:channel_slicing}.
Consider a convolutional layer with five channels whose deltas are $\Delta_{\text{ch}}^{(c_1)} =
1.9$, $\Delta_{\text{ch}}^{(c_2)} = 0.25$, $\Delta_{\text{ch}}^{(c_3)} = 2.1$,
$\Delta_{\text{ch}}^{(c_4)} = 0.25$, and $\Delta_{\text{ch}}^{(c_5)} = 0.5$, yielding a total
energy of $E_{\text{ch}} = 5.0$.
With a threshold of $\alpha = 0.8$, the algorithm sorts the channels by their delta in descending
order: $c_3, c_1, c_5, c_2, c_4$.
It then accumulates deltas until the threshold is reached. After including $c_3$ and $c_1$ ($2.1 +
1.9 = 4.0 \geq \alpha \cdot E_{\text{ch}} = 4.0$), channels $c_3$ and $c_1$ are retained and the
remaining three are excluded.

\vspace{0.5cm}
\begin{figure}[ht!]
    \centering
    \begin{tikzpicture}[
                fm/.style={draw, thick, minimum width=2.8cm, minimum height=2.8cm},
                kept/.style={fm, fill=orange!25, draw=black},
                excluded/.style={fm, fill=gray!10, draw=black, dashed},
        ]
    % Stacked channels (back to front: c5 back, c1 front), offset 0.6 each for more visible labels
    \node[excluded] (c5) at (2.4, 2.4) {};
    \node[font=\footnotesize] at (2.4, 3.5) {$\Delta_{c_5} = 0.5$};

    \node[excluded] (c4) at (1.8, 1.8) {};
    \node[font=\footnotesize] at (1.8, 2.9) {$\Delta_{c_4} = 0.25$};

    \node[kept] (c3) at (1.2, 1.2) {};
    \node[font=\footnotesize] at (1.2, 2.3) {$\Delta_{c_3} = 2.1$};

    \node[excluded] (c2) at (0.6, 0.6) {};
    \node[font=\footnotesize] at (0.6, 1.7) {$\Delta_{c_2} = 0.25$};

    \node[kept] (c1) at (0, 0) {};
    \node[font=\footnotesize] at (0, 1.1) {$\Delta_{c_1} = 1.9$};


    \end{tikzpicture}
    \caption{Channel-level slicing example for a convolutional layer with $|\mathcal{C}| = 5$ channels. With $\alpha = 0.8$ and $E_{\text{ch}} = 5.0$, channels $c_3$ and $c_1$ are retained (orange) and the remaining three are excluded (gray).}
    \label{fig:channel_slicing}
\end{figure}

Channels that are not retained have their contribution values set to zero.
This means they do not contribute to the backward analysis, which can have a side effect on
previous layers.
When the contribution of a channel is set to zero, neurons in earlier layers that would have
received contributions through this channel may end up with a lower or zero cumulative contribution
as well.
In the worst case, a neuron that is relevant to the prediction could be excluded from the slice
because its only contribution path passed through a filtered channel.
In our evaluation, we investigate this by testing different values of $\alpha$ and examining how
channel filtering affects slice quality.

Together, block slicing and channel slicing reduce the number of neurons considered during the
backward analysis.
First, entire blocks with low $\Delta_{\text{block}}$ are excluded.
Then, channels with low $\Delta_{\text{ch}}$ within the retained blocks are filtered out.
The next section addresses how the remaining neuron-level contributions are computed efficiently
using a vectorized backward analysis.

\section{Vectorized Backward Analysis}
\label{sec:vectorized_backward}

After block slicing and channel slicing have reduced the set of components to analyze,
the backward analysis computes neuron-level contributions for the remaining network using the
contribution rules described in Table~\ref{tab:operator_rules}.
The original NNSlicer implementation performs this computation by iterating over individual neurons
and synapses in nested loops.
For each output neuron $n$ in a layer, it visits every incoming synapse, computes the local
contribution, and accumulates the result.
While this implementation is correct, it becomes slow for residual networks with millions of
neurons and synapses.
Zhang et al.\ report that NNSlicer requires several minutes to compute the backward analysis for a
single slice for residual networks such as ResNet-18~\cite{zhang2020dynamic}.
This motivates our reimplementation of the backward analysis using vectorized tensor operations.


For our implementation, we use PyTorch, a widely used open-source deep learning
framework~\cite{paszke2019pytorch}.
In PyTorch, data is represented as \textit{tensors}, multi-dimensional arrays that generalize
scalars, vectors, and matrices to arbitrary dimensions:
\begin{itemize}
\item A \textit{scalar} is a 0-dimensional tensor holding a single value, e.g.\ $5$.
\item A \textit{vector} is a 1-dimensional tensor, e.g.\ $[1, 3, 7]$.
\item A \textit{matrix} is a 2-dimensional tensor with rows and columns, e.g.\ shape $(3 \times 4)$.
\item \textit{Higher-order tensors} generalize this to additional axes.
\end{itemize}
Tensors correspond directly to the objects used in linear algebra.
Vector and matrix operations such as element-wise multiplication, dot products, and matrix
multiplication are natively supported by PyTorch and follow standard linear algebra conventions.
An example of a higher-order tensor is the activation output of a convolutional layer with $C$
channels and spatial dimensions $H \times W$, which forms a 3-dimensional tensor of shape $(C
\times H \times W)$.
PyTorch provides many different operations to manipulate these tensors. Common operations include
element-wise addition, multiplication, and summation along a given dimension. These operations are
applied to all elements simultaneously, making them significantly faster than explicit loops.

Throughout this section, we use a concrete example of a fully connected layer to illustrate the
vectorized backward analysis in comparison to a loop-based implementation.
We consider a fully connected layer $l$ with $|\mathcal{N}_l| = m = 3$ output neurons and
$|\mathcal{N}_{l-1}| = n = 4$ input neurons.
By the definition of fully connected layers, every input neuron is connected to every output
neuron, giving $|\mathcal{S}_l| = m \cdot n = 12$ synapses.
For fully connected layers, the backward analysis computes the local contribution of each input
neuron using the weighted contribution rule:
\begin{equation}
    \textit{contrib}_i = \textit{CONTRIB}_n \cdot \Delta y \cdot w_i \cdot \Delta x_i
\end{equation}
where $\textit{CONTRIB}_n$ is the cumulative contribution of output neuron $n \in \mathcal{N}_l$
and $\Delta y$ is its activation delta.
Since we compute the local contribution of input neuron $i \in \mathcal{N}_{l-1}$ with respect to
$n$, $w_i$ denotes the weight of the synapse connecting $i$ to $n$, and $\Delta x_i$ is the
activation delta of $i$.

A possible loop-based implementation of this computation rule is shown in
Algorithm~\ref{alg:loop_backward}.
For each output neuron, we iterate over all input neurons and apply the contribution rule.
In our example with $m = 3$ output neurons and $n = 4$ input neurons, this results in $m \cdot n =
12$ iterations.
In general, this has a time complexity of $\mathcal{O}(m \cdot n)$. Since each iteration is
executed sequentially, this becomes expensive as the network grows.

\begin{algorithm}[ht]
\caption{Loop-based backward contribution for a fully connected layer}\label{alg:loop_backward}
\begin{algorithmic}[1]
\For{each output neuron $n \in \mathcal{N}_l$}
    \If{$\textit{CONTRIB}_n \neq 0$}
        \For{each input neuron $i \in \mathcal{N}_{l-1}$}
            \State $\textit{contrib}_i \gets \textit{CONTRIB}_n \cdot \Delta y \cdot w_i \cdot \Delta x_i$
        \EndFor
    \EndIf
\EndFor
\end{algorithmic}
\end{algorithm}

In our vectorized approach, we replace these loops with tensor operations, computing the
contributions for all connections in a single operation.
To make this concrete, Figure~\ref{fig:vec_inputs} shows the four input tensors for our previous
example.
The cumulative contributions $\textit{CONTRIB}_n$ of the three output neurons are stored as a
vector of size $m$.
The corresponding activation deltas $\Delta y$ are also stored as a vector of size $m$.
The synapses $\mathcal{S}_l$ connecting the $m$ output neurons to the $n$ input neurons are
represented as a weight matrix $W \in \mathbb{R}^{m \times n}$, where each row contains the weights
from one output neuron to all input neurons.
Finally, the input deltas $\Delta x$ form a vector of size $n$, one value per input neuron.

\begin{figure}[ht]
\centering
\begin{tikzpicture}[
        cell/.style={minimum width=1.2cm, minimum height=0.85cm, draw, font=\small, inner sep=0pt},
        lbl/.style={font=\normalsize\itshape},
        op/.style={font=\Large},
        shp/.style={font=\footnotesize, text=gray},
]

% CONTRIB_n vector (m=3)
\node[lbl] at (0, 3.8) {$\textit{CONTRIB}_n$};
\node[cell] at (0, 2.8) {1};
\node[cell] at (0, 1.95) {1};
\node[cell] at (0, 1.1) {$-1$};
\node[shp] at (0, 0.4) {$(m)$};

% dot
\node[op] at (1.1, 1.95) {$\cdot$};

% Delta y vector (m=3)
\node[lbl] at (2.2, 3.8) {$\Delta y$};
\node[cell] at (2.2, 2.8) {1};
\node[cell] at (2.2, 1.95) {2};
\node[cell] at (2.2, 1.1) {1};
\node[shp] at (2.2, 0.4) {$(m)$};

% dot
\node[op] at (3.3, 1.95) {$\cdot$};

% W matrix (m=3 x n=4)
\node[lbl] at (6.2, 3.8) {$W$};
\node[cell] at (4.4, 2.8) {3};
\node[cell] at (5.6, 2.8) {3};
\node[cell] at (6.8, 2.8) {$-2$};
\node[cell] at (8.0, 2.8) {$-1$};
\node[cell] at (4.4, 1.95) {3};
\node[cell] at (5.6, 1.95) {3};
\node[cell] at (6.8, 1.95) {$-3$};
\node[cell] at (8.0, 1.95) {$-1$};
\node[cell] at (4.4, 1.1) {2};
\node[cell] at (5.6, 1.1) {$-2$};
\node[cell] at (6.8, 1.1) {$-3$};
\node[cell] at (8.0, 1.1) {1};
\node[shp] at (6.2, 0.4) {$(m \times n)$};

% dot
\node[op] at (9.1, 1.95) {$\cdot$};

% Delta x vector (n=4)
\node[lbl] at (10.2, 3.8) {$\Delta x$};
\node[cell] at (10.2, 2.8) {1};
\node[cell] at (10.2, 1.95) {2};
\node[cell] at (10.2, 1.1) {1};
\node[cell] at (10.2, 0.25) {3};
\node[shp] at (10.2, -0.45) {$(n)$};

\end{tikzpicture}
\caption{Tensor inputs for the vectorized backward contribution of a fully connected layer with $m = 3$ output and $n = 4$ input neurons.}
\label{fig:vec_inputs}
\end{figure}

Since all inputs are stored as tensors, we can compute the contribution for all connections
simultaneously.
Broadcasting is a mechanism in PyTorch that automatically expands tensors along missing dimensions
so that element-wise operations can be applied between tensors of different shapes.
In our case, the vectors $\textit{CONTRIB}_n$ and $\Delta y$ of size $m$ are broadcast across the
$n$ columns of $W$, while $\Delta x$ of size $n$ is broadcast across its $m$ rows.
This allows us to compute the element-wise product $\textit{CONTRIB}_n \cdot \Delta y \cdot W \cdot
\Delta x$ in a single operation.
The result is a matrix of size $m \times n$, where each entry contains the local contribution value
for one connection.
Figure~\ref{fig:vec_contrib} shows this intermediate result. Each entry represents the local
contribution of one output-input connection. These are not yet the final neuron contributions but
serve as input for the next step.

\begin{figure}[ht]
\centering
\begin{tikzpicture}[
        cell/.style={minimum width=1.2cm, minimum height=0.85cm, draw, font=\small, inner sep=0pt},
        lbl/.style={font=\normalsize\itshape},
        shp/.style={font=\footnotesize, text=gray},
]

% contrib matrix (m=3 x n=4)
% Row 1: CONTRIB=1, Dy=1 -> 1*1*[3,3,-2,-1]*[1,2,1,3] = [3, 6, -2, -3]
% Row 2: CONTRIB=1, Dy=2 -> 1*2*[3,3,-3,-1]*[1,2,1,3] = [6, 12, -6, -6]
% Row 3: CONTRIB=-1, Dy=1 -> -1*1*[2,-2,-3,1]*[1,2,1,3] = [-2, 4, 3, -3]
\node[cell] at (0.6, 1.8) {3};
\node[cell] at (1.8, 1.8) {6};
\node[cell] at (3.0, 1.8) {$-2$};
\node[cell] at (4.2, 1.8) {$-3$};
\node[cell] at (0.6, 0.95) {6};
\node[cell] at (1.8, 0.95) {12};
\node[cell] at (3.0, 0.95) {$-6$};
\node[cell] at (4.2, 0.95) {$-6$};
\node[cell] at (0.6, 0.1) {$-2$};
\node[cell] at (1.8, 0.1) {4};
\node[cell] at (3.0, 0.1) {3};
\node[cell] at (4.2, 0.1) {$-3$};
\node[shp] at (2.4, -0.6) {$(m \times n)$};

\end{tikzpicture}
\caption{Intermediate contribution matrix after element-wise multiplication with broadcasting. Each entry represents the raw contribution value for one output-input connection.}
\label{fig:vec_contrib}
\end{figure}

NNSlicer does not use the raw contribution values. Instead, it accumulates only the sign of each
local contribution:
\begin{equation}
    \textit{CONTRIB}_{n_i} \mathrel{+}= \text{sign}(\textit{contrib}_i)
\end{equation}
We apply the sign function element-wise to the result matrix, replacing each entry with $+1$, $-1$,
or $0$.
The resulting sign matrix is shown in Figure~\ref{fig:vec_sign}, where each entry is now $+1$ or
$-1$, indicating whether the connection has a positive or negative contribution.

\begin{figure}[ht]
\centering
\begin{tikzpicture}[
        cell/.style={minimum width=1.2cm, minimum height=0.85cm, draw, font=\small, inner sep=0pt},
        lbl/.style={font=\normalsize\itshape},
        op/.style={font=\Large},
        shp/.style={font=\footnotesize, text=gray},
        arrow/.style={->, thick, >=stealth},
]

% sign matrix (m=3 x n=4)
\node[cell] at (0.6, 1.8) {$+1$};
\node[cell] at (1.8, 1.8) {$+1$};
\node[cell] at (3.0, 1.8) {$-1$};
\node[cell] at (4.2, 1.8) {$-1$};
\node[cell] at (0.6, 0.95) {$+1$};
\node[cell] at (1.8, 0.95) {$+1$};
\node[cell] at (3.0, 0.95) {$-1$};
\node[cell] at (4.2, 0.95) {$-1$};
\node[cell] at (0.6, 0.1) {$-1$};
\node[cell] at (1.8, 0.1) {$+1$};
\node[cell] at (3.0, 0.1) {$+1$};
\node[cell] at (4.2, 0.1) {$-1$};
\node[shp] at (2.4, -0.6) {$(m \times n)$};

\end{tikzpicture}
\caption{After applying sign, each entry is $+1$, $0$, or $-1$.}
\label{fig:vec_sign}
\end{figure}

Summing over the rows yields the contribution vector $[1, 3, -1, -3]$ for the four input neurons.
This vector becomes the new $\textit{CONTRIB}$ input for the next layer, and the process repeats
until the first layer of the network is reached.

The example above illustrates the vectorized backward analysis for a fully connected layer.
The same principle applies to convolutional layers, where the computation is performed over local
input patches and kernel weights.
These patches and kernel weights are represented as tensors, allowing the same element-wise
operations to be used.
In general, we use vectorized tensor operations throughout all phases of our implementation.
While we evaluate our approach on a CPU, tensor operations can also be executed on a GPU, which
could further improve performance due to its ability to run thousands of operations in parallel.

\vspace{0.2cm}

In this chapter, we presented a four-step approach that extends the original workflow of NNSlicer. 
Delta aggregation provides an input-specific indicator derived from the forward analysis. 
Block slicing and channel slicing use these aggregated deltas to reduce the scope of the backward
analysis by skipping irrelevant blocks and channels. 
Finally, the vectorized backward analysis replaces the original loop-based computation with tensor
operations. 
Together, these steps aim to make neural network slicing more efficient while preserving the
quality of the resulting slices.
In the following chapter, we evaluate this approach by measuring both the computational speedup and
the impact on slice quality.



\chapter{Evaluation}
\label{chap:evaluation}

In this chapter, we evaluate our reimplementation of NNSlicer and the proposed block-based
extensions.
We assess runtime efficiency, slice quality, and interpretability by the following three research
questions:

\begin{itemize}
    \item \textbf{RQ1:} How does the block-based approach reduce runtime and slice size compared to the neuron-level approach?
    \item \textbf{RQ2:} Do slices produced by the block-based approach retain information relevant for classification?
    \item \textbf{RQ3:} Do slicing results reveal interpretable class-specific patterns at channel and block level?
\end{itemize}

The chapter is structured as follows: First, we describe the experimental setup, including the
models and hardware used.
We then present the results for each research question, followed by a discussion of threats to
validity.

\section{Experimental Setup}

This section describes the experimental setup used in our evaluation, including the datasets,
models, and hardware.

\paragraph{Datasets.}
For our evaluation we use two image classification datasets: CIFAR-10~\cite{krizhevsky2009learning}
and ImageNet~\cite{deng2009imagenet, russakovsky2015imagenet}.
CIFAR-10 contains 60{,}000 color images of size $32 \times 32$ pixels, divided into 10 classes with
6{,}000 images each.
The classes cover six animals (bird, cat, deer, dog, frog, horse) and four vehicles (airplane,
automobile, ship, truck).
ImageNet is a large-scale dataset with over 1.2 million training images across 1{,}000 classes.
These classes cover different objects, animals, and scenes.


\paragraph{Models.}
We use two types of ResNet models in our evaluation.
On the one hand, we use three models from the official ResNet family: ResNet-18, ResNet-34, and
ResNet-101.
They are particularly interesting for RQ1 as they grow in size, allowing us to study how slicing
scales to larger networks.
To work with these models, we use the pretrained weights provided by the torchvision
library~\cite{paszke2019pytorch}.

On the other hand, we further created a custom ResNet model trained on the CIFAR-10 dataset
similarly to Zhang et al.~\cite{zhang2020dynamic}.
The model uses three groups of BasicBlock residual blocks with 16, 32, and 64 channels respectively
and follows the CIFAR ResNet architecture proposed by He et al.~\cite{he2016deep}. 
We chose 27 blocks because a higher number of blocks provides more granularity for evaluating
block-level slicing and makes the effect of different $\beta$ thresholds more visible.
Additionally, CIFAR-10 is bundled with torchvision and can be downloaded automatically during
training, making experiments easily reproducible.
The CIFAR-10 resnet model has been trained manually for 150 epochs using SGD with a learning rate
of 0.1 and momentum of 0.9.
A test accuracy of 91.08\% is achieved across all ten classes. All models used in our evaluation
are summarized in Table~\ref{tab:model_parameters}.

\paragraph{Precomputed Profiling.}
Slicing in NNSlicer consists of three phases: profiling, forward analysis, and backward analysis.
The profiling phase is a one-time preprocessing step that computes the mean activation of each
neuron over a dataset.
For the ImageNet models, profiling was performed over the ImageNet Large Scale Visual Recognition
Challenge 2012 (ILSVRC2012) validation set~\cite{deng2009imagenet, russakovsky2015imagenet}.
The dataset can be downloaded from the official ImageNet website and has a size of
6.3\,GB~\cite{imagenet2012}.
For ResNet-CIFAR, profiling was performed over the CIFAR-10 training set.
Since profiling only needs to be computed once per model, we precompute profiling results and do
not include them in our runtime evaluation.
We focus exclusively on backward analysis runtimes, as this is the computationally expensive step
that must be repeated for every target class and input image.

\paragraph{Parameters.}
Our approach uses three threshold parameters.
The first, $\theta$, originates from NNSlicer and controls the neuron-level filtering during the
backward analysis.
The original paper does not specify a value for $\theta$.
We inspected the NNSlicer source code and found a corresponding threshold value of $0.2$.
We validated this choice by reproducing the example network from the original NNSlicer paper and
confirming that our implementation produces the same slice with $\theta = 0.2$.
We keep $\theta$ fixed at this value throughout all experiments.
The remaining two parameters, $\alpha$ and $\beta$, control our block and channel filtering
respectively.
Unlike $\theta$, these are varied across experiments to study their effect on runtime and slice
quality.

\paragraph{Hardware.}
All forward and backward analysis experiments in RQ1--RQ3 were conducted on an Apple MacBook Pro
with an M2 Pro chip and 16\,GB of RAM.
Profiling requires a forward pass for every image in the dataset, which is time-consuming for large
datasets like ImageNet.
Since the MacBook does not have a dedicated GPU, profiling was performed on Google Colab using an
NVIDIA A100 GPU~\cite{bisong2019colab}.

\begin{table}[ht]
\centering
\caption{Models used in the evaluation with their parameter counts and residual blocks. ImageNet models use pretrained torchvision weights. ResNet-CIFAR is trained on CIFAR-10}
\label{tab:model_parameters}
\begin{tabular}{l l r c}
\toprule
\textbf{Model} & \textbf{Dataset} & \textbf{\# Parameters} & \textbf{\# Residual Blocks} \\
\midrule
ResNet-18   & ImageNet & 11.7M   & 8  \\
ResNet-34   & ImageNet & 21.8M   & 16 \\
ResNet-101  & ImageNet & 44.5M   & 33 \\
\midrule
ResNet-CIFAR & CIFAR-10 & 855{,}770 & 27 \\
\bottomrule
\end{tabular}
\end{table}

\noindent

With the experimental setup defined, we continue with the results of our evaluation.

\section{Results}

We present results for each research question below.
For RQ1, we compare runtime and slice sizes across different configurations.
For RQ2, we assess slice quality through model pruning and counterfactual reasoning.
RQ3 examines whether slicing reveals interpretable class-specific patterns.

\subsection*{RQ1: How does the block-based approach reduce runtime and slice size compared to the neuron-level approach?}

In RQ1 we evaluate the runtime and slice size of generating a single slice.
We use the ImageNet ResNet models to study how slicing scales with increasing model depth and size.
Additionally, we use ResNet-CIFAR to provide an isolated comparison between our vectorized
implementation and a loop-based variant on the same hardware.

We report slice sizes in terms of neurons and synapses.
A neuron is considered part of the slice if its contribution value is non-zero.
For synapses, we count the incoming connections to each active output neuron: if an output neuron
is in the slice, all its incoming connections from the previous layer are counted as active.
Both metrics are computed only over layers with learnable weights, i.e., convolutional and fully
connected layers.

For the runtime, we measure purely the backward analysis time, as done in the original NNSlicer
paper.
Each reported runtime is the median of five runs for the same target class and input image to
account for system-level variance.
To verify that this is sufficient, we repeated slicing across 20 randomly selected target classes
for ResNet-18 and ResNet-34.
The variance across classes is negligible for neuron and synapse ratios as well as runtime.
This is expected, as the backward analysis operates on tensors whose dimensions are determined by
the model architecture, not by the target class.

In this section, we first create a baseline by running backward analysis without channel or block
slicing.
This baseline serves as a reference for comparing against the original NNSlicer and for measuring
the speedup from our block-based extensions.
Table~\ref{tab:rq1-baseline} reports the backward pass time per slice with our implementation.

\begin{table}[ht]
\centering
\caption{Backward analysis time per slice without channel or block filtering (baseline). The upper section isolates the effect of vectorization on CIFAR-10 ResNet. The lower section shows ImageNet ResNet models with our vectorized implementation.}
\label{tab:rq1-baseline}
\begin{tabular*}{\textwidth}{l @{\extracolsep{\fill}} r r r}
\toprule
\textbf{Model} & \textbf{Neurons (\%)} & \textbf{Synapses (\%)} & \textbf{Time / slice} \\
\midrule
ResNet-CIFAR (vectorized) & 47.0 & 41.2 & 0.15\,s \\
ResNet-CIFAR (loop)       & 47.0 & 41.2 & 639.77\,s \\
\midrule
ResNet-18    & 36.2 & 38.2 & 2.00\,s \\
ResNet-34    & 41.0 & 41.2 & 3.48\,s \\
ResNet-101   & 41.8 & 34.0 & 7.67\,s \\
\bottomrule
\end{tabular*}
\end{table}

The upper section of the table isolates the effect of our vectorized implementation by comparing it
against a loop-based variant on the same hardware.
Both produce identical slices, but the vectorized version completes in 0.15\,s compared to
639.77\,s for the loop-based variant. This confirms the performance gain of replacing
neuron-by-neuron loops with tensor operations.

The lower section shows that all ImageNet models can be sliced within seconds with our
implementation.
ResNet-18 completes a single slice in about 2\,s. Even the largest model ResNet-101 with 44.5M
parameters is sliced in about 7.7 seconds.
Slices typically retain 36--47\,\% of neurons and 34--41\,\% of synapses, showing that roughly half
of the network is identified as irrelevant without block or channel slicing.

As a reference point, the original NNSlicer reports a runtime of approximately 543\,s for a single
slice of ResNet-18.
However, a direct comparison is not possible as this number was obtained on different hardware and
a TensorFlow-based implementation.
The original authors also propose processing multiple slices in parallel, which reduces the
backward pass to approximately 40\,s for ResNet-18.
Our single-slice implementation already outperforms this, and a similar batching strategy could be
explored in future work.

\begin{table}[ht]
\centering
\caption{Effect of block slicing on runtime and slice size for three $\beta$ thresholds. Skipped blocks are excluded from backward analysis entirely}
\label{tab:rq1-block}
\begin{tabular*}{\textwidth}{l @{\extracolsep{\fill}} c r r r c}
\toprule
\textbf{Model} & $\boldsymbol{\beta}$ & \textbf{Time} & \textbf{Neurons (\%)} & \textbf{Synapses (\%)} & \textbf{Skipped} \\
\midrule
\multirow{3}{*}{ResNet-18}
        & 0.80 & 1.62\,s & 32.2 & 28.7 & 2/8 \\
        & 0.50 & 1.28\,s & 21.9 & 16.5 & 4/8 \\
        & 0.20 & 0.86\,s & 0.9 & 3.4 & 7/8 \\
\midrule
\multirow{3}{*}{ResNet-34}
        & 0.80 & 2.77\,s & 37.3 & 32.6 & 4/16 \\
        & 0.50 & 1.84\,s & 21.8 & 14.5 & 10/16 \\
        & 0.20 & 1.22\,s & 1.1 & 4.0 & 14/16 \\
\midrule
\multirow{3}{*}{ResNet-101}
        & 0.80 & 6.50\,s & 35.8 & 27.0 & 7/33 \\
        & 0.50 & 4.02\,s & 27.1 & 16.9 & 18/33 \\
        & 0.20 & 2.03\,s & 13.1 & 6.0 & 28/33 \\
\bottomrule
\end{tabular*}
\end{table}

With the baseline established, we can evaluate whether channel and block slicing can further reduce
runtime or slice size.
At block slicing, entire residual blocks whose energy falls below a threshold~$\beta$ are skipped.
We evaluate three different values for $\beta$.
A conservative filter with $0.80$ that retains most of the blocks. A balanced middle ground with
$\beta = 0.50$.
The third one is aggressive filter with $\beta = 0.20$ that retains only the blocks with strongest
deltas.
From Table~\ref{tab:rq1-block} we can observe that block slicing reduces both slice size and
runtime across all models.
Lowering~$\beta$ consistently increases the number of skipped blocks.
These skipped blocks are excluded from the backward analysis entirely and their neurons and
synapses no longer appear in the slice.
As a result, both slice size and runtime decrease with fewer processed blocks.
For ResNet-101, 7 of 33 blocks are skipped for $\beta = 0.80$. The runtime drops from 7.67\,s to
6.50\,s and the neuron ratio decreases from 41.8\,\% to 35.8\,\%.
At $\beta = 0.20$, 28 of 33 blocks are skipped. This further reduces runtime to 2.00\,s and the
neuron ratio to 13.1\,\%.
The same trend holds for the other two models ResNet-18 and ResNet-34.
This confirms that block slicing with lower values for $\beta$ creates smaller and faster slices.

\begin{table}[ht]
\centering
\caption{Effect of channel slicing on runtime and slice size for three $\alpha$ thresholds. Lower $\alpha$ retains fewer channels, reducing slice size only}
\label{tab:rq1-channel}
\begin{tabular*}{\textwidth}{l @{\extracolsep{\fill}} c r r r}
\toprule
\textbf{Model} & $\boldsymbol{\alpha}$ & \textbf{Time} & \textbf{Neurons (\%)} & \textbf{Synapses (\%)} \\
\midrule
\multirow{3}{*}{ResNet-18}
        & 0.80 & 2.03\,s & 25.0 & 27.6 \\
        & 0.50 & 1.97\,s & 13.4 & 15.4 \\
        & 0.20 & 2.05\,s & 4.6 & 5.5 \\
\midrule
\multirow{3}{*}{ResNet-34}
        & 0.80 & 3.51\,s & 29.0 & 29.9 \\
        & 0.50 & 3.44\,s & 15.4 & 16.6 \\
        & 0.20 & 3.53\,s & 5.3 & 5.9 \\
\midrule
\multirow{3}{*}{ResNet-101}
        & 0.80 & 7.71\,s & 27.4 & 23.7 \\
        & 0.50 & 7.62\,s & 14.8 & 13.6 \\
        & 0.20 & 7.69\,s & 5.1 & 5.0 \\
\bottomrule
\end{tabular*}
\end{table}

While block slicing operates at the coarsest level of blocks, channel slicing removes channels
whose energy contribution falls below a threshold~$\alpha$.
These filtered channels are excluded from the backward analysis as their contribution values within
the tensors are set to zero.
Similar to block slicing, we evaluate three different values for $\alpha$: $0.80$, $0.50$, and
$0.20$.
The results in Table~\ref{tab:rq1-channel} show that runtime remains nearly identical to the
baseline across all $\alpha$ values.
This is a direct consequence of our implementation: filtered channels have their contribution
values set to zero within the convolutional tensors, but the tensor dimensions remain unchanged.
The underlying tensor operations therefore stay the same size and do not affect the computational
cost.

However, while channel slicing does not provide a speedup, it creates smaller slices.
Filtered channels reduce the number of active neurons in each convolutional layer, which in turn
reduces the synapse count.
For example, at $\alpha = 0.80$ on ResNet-18 the neuron percentage drops from 36.2\,\% to 25.0\,\%
and the synapse percentage from 38.2\,\% to 27.6\,\%.
Similar to block slicing, lowering $\alpha$ produces smaller slices.
With a low value of $\alpha = 0.20$, only 5.1\,\% of neurons and 5.0\,\% of synapses remain on
ResNet-101.

Finally, we evaluate the combined effect of both strategies.
Block slicing first skips entire blocks, then channel slicing filters within the remaining blocks.
This should yield both smaller slices and lower runtime.
We evaluate this for a single combination of $\alpha = 0.85$ and $\beta = 0.80$.
This combination preserves a high slice quality and demonstrates a meaningful choice of thresholds,
which we validate in RQ2.

\begin{table}[ht]
\centering
\caption{Combined channel and block slicing with $\alpha = 0.85$ and $\beta = 0.80$, showing the cumulative effect of both slicing strategies}
\label{tab:rq1-combined}
\begin{tabular*}{\textwidth}{l @{\extracolsep{\fill}} r r r c}
\toprule
\textbf{Model} & \textbf{Time} & \textbf{Neurons (\%)} & \textbf{Synapses (\%)} & \textbf{Skipped} \\
\midrule
ResNet-18  & 1.69\,s & 24.1 & 22.2 & 2/8 \\
ResNet-34  & 2.95\,s & 28.6 & 25.6 & 4/16 \\
ResNet-101 & 6.20\,s & 25.9 & 20.5 & 7/33 \\
\bottomrule
\end{tabular*}
\end{table}

Table~\ref{tab:rq1-combined} shows that combining both techniques yields lower neuron and synapse
percentages than either technique alone.
For example, on ResNet-18 the combined configuration reduces neuron percentages from 36.2\,\% to
24.1\,\% and synapse percentages from 38.2\,\% to 22.2\,\% compared to the baseline slice.
The backward pass time is also reduced from 2.00\,s to 1.69\,s, since block skipping eliminates
entire residual blocks from the backward pass.
The same pattern holds for ResNet-34 and ResNet-101.
This confirms that the two strategies are complementary: block slicing contributes runtime savings
by skipping entire blocks, while channel slicing contributes slice precision by filtering within
the remaining blocks.

\vspace{0.2cm}

In summary, our implementation reduces the backward analysis for ResNet-18 from over 540\,s to
about 2\,s.
It is now also possible to slice large residual networks like ResNet-101 in seconds.
Channel slicing reduces slice size, while block slicing reduces both runtime and slice size.
A combined configuration ($\alpha{=}0.85$, $\beta{=}0.80$) reduces neuron and synapse percentages
to roughly half of the baseline values.
However, smaller slices risk losing relevant neurons.
This raises the question of whether the retained slices still capture the information needed for
classification, which we address in RQ2.

\subsection*{RQ2: Do slices produced by the block-based approach retain information relevant for classification?}

We address this question through two strategies.
First, we use model pruning to remove components outside the slice and measure the effect on
prediction accuracy.
Second, we apply counterfactual reasoning to test whether targeted changes to slice and non-slice
regions affect the prediction.

\subsection*{Model Pruning}

Model pruning is a technique that reduces the size of a neural network by removing parameters that
contribute little to its predictions.
The idea was first introduced by LeCun et al.~\cite{lecun1989optimal}, who showed that removing
low-importance weights from a trained network can reduce its size with minimal loss in accuracy.

Our pruning follows the hierarchical structure of the slicing approach and operates at three
levels.
First, entire residual blocks whose contributions fall below the threshold $\beta$ are removed by
zeroing all their weights.
Next, channels that have been filtered out by $\alpha$ are not considered in the slice and pruned
out.
Finally, on the neuron level, we prune synapses that have no contribution in the slice.
This produces a pruned model that contains only the subnetwork identified by the slice.

We evaluate pruning on our custom residual network trained on CIFAR-10.
The model has 27 residual blocks and is deep enough to be meaningful, but small enough to make
pruning results reproducible.
For each class, we compute individual slices $\mathcal{M}_{C,1}, \dots, \mathcal{M}_{C,10}$ from
ten training images and aggregate them via $\mathcal{M}_C = \bigcup_{i=1}^{10} \mathcal{M}_{C,i}$.
This produces a class-specific slice, e.g. the union of ten different cat image slices for the cat
class.
The pruned model is then evaluated on the full CIFAR-10 test set of 10{,}000 images using the
following metrics:
\begin{itemize}
    \item \textbf{Target Accuracy}: average accuracy on the class the slice was computed for. A high value indicates that the slice retains the relevant information.
    \item \textbf{Non-Target Accuracy}: average accuracy on the remaining nine classes. A low value indicates that the pruned model is specific to the target class.
    \item \textbf{Model Size}: the percentage of non-zero parameters remaining after pruning.
\end{itemize}

\noindent
Similar to RQ1 we evaluate three configurations: channel-only pruning with varying $\alpha$.
Block-only pruning with varying $\beta$. A combined configuration where $\alpha$ is fixed at $0.85$
while $\beta$ varies.
In each table, the first row shows the baseline, which corresponds to the original NNSlicer
approach without block or channel filtering.

\begin{table}[ht]
\centering
\caption{Channel pruning results for varying $\alpha$. Target and non-target accuracy after zeroing out channels excluded by the slice}
\label{tab:channel_pruning}
\begin{tabular}{cccc}
\toprule
\textbf{$\alpha$} & \textbf{Target} & \textbf{Non-Target} & \textbf{Model Size} \\
\midrule
-- & 93.6 & 74.5 & 53.1\% \\
\midrule
0.95 & 92.3 & 71.8 & 51.9\% \\
0.90 & 91.7 & 71.0 & 50.6\% \\
0.85 & 91.2 & 70.4 & 49.1\% \\
0.80 & 70.5 & 57.2 & 47.6\% \\
0.75 & 64.8 & 42.9 & 46.2\% \\
0.70 & 50.7 & 26.1 & 44.5\% \\
\bottomrule
\end{tabular}
\end{table}

For channel-only pruning, Table~\ref{tab:channel_pruning} shows that filters of $\alpha = 0.95$
down to $\alpha = 0.85$ retain high target accuracy above 91\%.
The non-target accuracy drops slightly from 74.5\% to 70.4\%, showing that the pruned model becomes
more focused on the target class.
However, target accuracy drops sharply at $\alpha = 0.80$, falling to 70.5\%.
This steep decline can be explained by the nature of channel pruning.
Each channel in a convolutional layer extracts a specific feature and when a channel is removed,
that feature is lost entirely.
Even channels with relatively low energy may carry features that are critical for classification.
This makes channel pruning sensitive to the threshold and suggests that only the least active
channels can be safely removed.

Block pruning shows a different picture in Table~\ref{tab:block_pruning}.
At $\beta = 0.8$, the model is reduced to 39.7\% of its original size by skipping 7 of 27 blocks.
Target accuracy remains high at 91.9\%, close to the baseline of 93.6\%.
Non-target accuracy drops from 74.5\% to 53.7\%, meaning the pruned model becomes more
class-specific.
\begin{table}[ht]
\centering
\caption{Block pruning results for varying $\beta$. Target and non-target accuracy after zeroing out neurons in blocks excluded by the slice}
\label{tab:block_pruning}
\begin{tabular}{ccccc}
\toprule
\textbf{$\beta$} & \textbf{Target} & \textbf{Non-Target} & \textbf{Block Skip} & \textbf{Model Size} \\
\midrule
-- & 93.6 & 74.5 & 0/27 & 53.1\% \\
\midrule
0.90 & 92.7 & 66.9 & 3/27 & 48.8\% \\
0.80 & 91.9 & 53.7 & 7/27 & 39.7\% \\
0.70 & 79.6 & 33.1 & 10/27 & 26.2\% \\
0.60 & 65.5 & 24.3 & 13/27 & 18.9\% \\
0.50 & 62.7 & 22.9 & 14/27 & 16.6\% \\
0.20 & 27.6 & 12.3 & 21/27 & 9.7\% \\
\bottomrule
\end{tabular}
\end{table}

This works because the skip connections in the residual architecture allow entire blocks to be
removed without disrupting the signal flow.
At $\beta = 0.5$, 14 of 27 blocks are skipped and only 16.6\% of parameters remain.
Target accuracy drops to 62.7\% with a non-target accuracy of 22.9\%.
At $\beta = 0.2$, 21 of 27 blocks are skipped and target accuracy drops to 27.6\%.
At this point, the filtering becomes too aggressive and critical parts of the network are lost.

\begin{table}[ht]
\centering
\caption{Combined pruning results with fixed $\alpha = 0.85$ and varying $\beta$, showing how block filtering complements channel filtering}
\label{tab:combined_pruning}
\begin{tabular}{ccccc}
\toprule
\textbf{$\beta$} & \textbf{Target} & \textbf{Non-Target} & \textbf{Block Skip} & \textbf{Model Size} \\
\midrule
-- & 93.6 & 74.5 & 0/27 & 53.1\% \\
\midrule
0.90 & 91.4 & 59.9 & 3/27 & 44.5\% \\
0.80 & 89.0 & 48.1 & 7/27 & 36.4\% \\
0.70 & 71.0 & 31.8 & 10/27 & 23.9\% \\
0.60 & 63.2 & 24.9 & 13/27 & 17.1\% \\
\bottomrule
\end{tabular}
\end{table}

Since channel pruning is sensitive to the threshold, we fix $\alpha = 0.85$ as a conservative value
that retains high target accuracy.
Block pruning allows for more aggressive filtering without losing critical features, so we vary
$\beta$.
Combining both strategies further reduces model size (Table~\ref{tab:combined_pruning}).
At $\beta = 0.9$, the combined configuration achieves 91.4\% target accuracy at 44.5\% model size,
compared to 48.8\% with block pruning only.
The non-target accuracy drops to 59.9\%, compared to 66.9\% with block pruning alone, showing that
the combined approach produces a more focused subnetwork.
At $\beta = 0.7$, the model retains 23.9\% of its parameters with 71.0\% target accuracy and a
non-target accuracy of 31.8\%.
Block-level pruning provides the largest gains, while channel filtering contributes additional size
reductions and improves class specificity.

Figure~\ref{fig:pruning_tradeoff} shows the relationship between threshold value and target
accuracy for all three strategies.
Block pruning maintains high target accuracy even at aggressive thresholds, while channel pruning
degrades rapidly below $\alpha = 0.85$.

\begin{figure}[ht]
\centering
\begin{tikzpicture}
\begin{axis}[
        width=0.85\textwidth,
        height=0.50\textwidth,
        xlabel={Threshold},
        ylabel={Target Accuracy (\%)},
        xmin=0.2, xmax=1.0,
        x dir=reverse,
        ymin=0, ymax=100,
        xtick={1.0, 0.9, 0.8, 0.7, 0.6, 0.5, 0.4, 0.3},
        grid=major,
        grid style={gray!30},
        thick,
        legend pos=south west, legend style={font=\footnotesize},
]
\addplot[color=blue, mark=square*, mark size=2, line width=1.2pt] coordinates {
        (1.0, 93.6) (0.90, 91.7) (0.80, 70.5) (0.70, 50.7) (0.60, 29.5) (0.50, 19.1) (0.40, 36.5)
    (0.30, 25.9) (0.20, 6.0)
};
\addlegendentry{Channel ($\alpha$)}
\addplot[color=red, mark=triangle*, mark size=2, line width=1.2pt] coordinates {
        (1.0, 93.6) (0.90, 92.7) (0.80, 91.9) (0.70, 79.6) (0.60, 65.5) (0.50, 62.7) (0.40, 48.0)
    (0.30, 43.6) (0.20, 27.6)
};
\addlegendentry{Block ($\beta$)}
\addplot[color=gray, mark=diamond*, mark size=2, line width=1.2pt, dashed] coordinates {
        (1.0, 93.6) (0.90, 91.4) (0.80, 89.0) (0.70, 71.0) (0.60, 63.2) (0.50, 62.6) (0.40, 48.5)
    (0.30, 40.8) (0.20, 28.4)
};
\addlegendentry{Combined ($\alpha{=}0.85$, $\beta$)}
\end{axis}
\end{tikzpicture}
\caption{Target accuracy as a function of the threshold for channel, block, and combined pruning. Block slicing retains high accuracy even at aggressive thresholds, while channel slicing degrades earlier}
\label{fig:pruning_tradeoff}
\end{figure}

Overall, the pruning results confirm that the slices produced by our approach capture the relevant
subnetwork for each class.
Channel filtering alone degrades accuracy quickly because removing a channel sets its contributions
to zero, which can also cause neurons in subsequent layers to lose their contributions.
Block-level filtering may work better because even when a block is skipped, its contribution can
still flow through the residual skip connection.
Combining both strategies yields similar accuracy to block pruning alone but achieves additional
size reductions.
In the following section, we evaluate slice quality from a different angle using counterfactual
reasoning.

\subsection*{Counterfactual Reasoning}

Counterfactual reasoning asks what would have happened if conditions had been
different~\cite{pearl2009causality}.
We apply this idea to evaluate slice quality by posing two complementary what-if questions:

\begin{itemize}
    \item \textbf{Necessity}: What happens if we progressively remove neurons \emph{inside} the slice? If the slice correctly identifies the relevant subnetwork, removing its neurons should rapidly degrade the prediction of the model for the target class.
    \item \textbf{Sufficiency}: What happens if we progressively remove neurons \emph{outside} the slice? If the slice is sufficient, the prediction should remain stable even as non-slice neurons are removed, because the essential computation is preserved.
\end{itemize}

Given a slice $\mathcal{M}_C = (\mathcal{N}_C, \mathcal{S}_C)$, the necessity test scales down the
activations of all neurons $n \in \mathcal{N}_C$ by a factor $s$ ranging from 1.0 to 0.5. At $s =
1.0$, the activations remain unchanged. As $s$ decreases toward 0.5, the activations are
progressively halved, simulating the gradual removal of slice neurons.
The sufficiency test applies the same scaling to all neurons $n \notin \mathcal{N}$, leaving only
the slice neurons at full strength.
In both cases, we apply softmax to the model output and measure the average target class
probability across all ten CIFAR-10 classes using ten test images per class.
We only evaluate images that the model classifies correctly, ensuring that the baseline prediction
is meaningful.

We compare three slice configurations to understand how filtering affects slice quality:

\begin{itemize}
    \item \textbf{Baseline} ($\alpha{=}\text{None},\;\beta{=}\text{None}$): the original NNSlicer without any channel or block filtering, producing the largest and most inclusive slices.
    \item \textbf{Combined} ($\alpha{=}0.85,\;\beta{=}0.8$): the best combined configuration from the pruning evaluation, which applies moderate filtering while preserving target accuracy.
    \item \textbf{Aggressive} ($\alpha{=}0.5,\;\beta{=}0.5$): a deliberately over-filtered configuration that removes a large portion of channels and blocks, producing much smaller but potentially incomplete slices.
\end{itemize}

The aggressive configuration serves as a negative control. An over-filtered slice should show
weaker necessity and sufficiency, as important neurons are missing from the slice.

Figure~\ref{fig:counterfactual_necessity} presents the necessity results.
The baseline drops fastest, with the combined configuration close behind.
At 15\,\% reduction, target probability falls to 55\,\% for the baseline and 71\,\% for the
combined configuration.
By 25\,\%, both are below 22\,\%.
The rapid drop confirms that the neurons identified by the slice are critical for the prediction.
Even a small reduction in their activations significantly degrades the output.
The aggressive configuration drops noticeably slower, still retaining 77\,\% target probability at
20\,\% reduction.
Because aggressive filtering has already excluded many important neurons from the slice, the
remaining slice neurons are less concentrated in their relevance. Scaling them down therefore has a
weaker effect on the prediction.

\begin{figure}[ht]
\centering
\begin{tikzpicture}
\begin{axis}[
        width=0.85\textwidth,
        height=0.50\textwidth,
        xlabel={Slice reduction (\%)},
        ylabel={Target probability (\%)},
        xmin=-2, xmax=52,
        ymin=0, ymax=110,
        xtick={0, 10, 20, 30, 40, 50},
        ytick={0, 20, 40, 60, 80, 100},
        grid=major,
        grid style={gray!30},
        thick,
        legend pos=north east, legend style={font=\footnotesize},
]
\addplot[color=red, mark=triangle*, mark size=2, line width=1.2pt] coordinates {
        (0, 98.2) (5, 94.7) (10, 79.9) (15, 54.6)
        (20, 24.4) (25, 11.4) (30, 9.3) (35, 8.6)
        (40, 8.0) (45, 7.5) (50, 7.0)
};
\addlegendentry{Baseline}
\addplot[color=blue, mark=square*, mark size=2, line width=1.2pt] coordinates {
        (0, 98.2) (5, 96.2) (10, 88.7) (15, 71.1)
        (20, 48.1) (25, 23.8) (30, 12.1) (35, 8.3)
        (40, 7.4) (45, 6.8) (50, 6.3)
};
\addlegendentry{$\alpha{=}0.85,\;\beta{=}0.8$}
\addplot[color=gray, mark=diamond*, mark size=2, line width=1.2pt, dashed] coordinates {
        (0, 98.2) (5, 97.0) (10, 93.9) (15, 87.7)
        (20, 76.9) (25, 62.5) (30, 47.3) (35, 32.7)
        (40, 21.4) (45, 14.0) (50, 10.3)
};
\addlegendentry{$\alpha{=}0.5,\;\beta{=}0.5$}
\end{axis}
\end{tikzpicture}
\caption{Necessity evaluation on ResNet-CIFAR. Target probability drops as neurons inside the slice are progressively scaled down.}
\label{fig:counterfactual_necessity}
\end{figure}

The sufficiency results are shown in Figure~\ref{fig:counterfactual_sufficiency}.
The baseline and combined configurations remain stable up to approximately 30\,\% reduction,
staying above 90\,\% target probability.
This indicates that the neurons outside the slice contribute little to the prediction, so removing
them has minimal effect.
Beyond 30\,\%, the prediction begins to degrade as more neurons are removed that, while not in the
slice, still play a supporting role.
At 50\,\% reduction, they decline to 37.3\,\% and 33.3\,\%.
The aggressive configuration drops earlier and reaches only 22.8\,\% at 50\,\% reduction.
Because aggressive filtering already excluded important neurons from the slice, these neurons are
now among those being kept, and their removal in the sufficiency test is less damaging. Instead,
the missing slice neurons cause the earlier decline.

\begin{figure}[ht]
\centering
\begin{tikzpicture}
\begin{axis}[
        width=0.85\textwidth,
        height=0.50\textwidth,
        xlabel={Non-slice reduction (\%)},
        ylabel={Target probability (\%)},
        xmin=-2, xmax=52,
        ymin=0, ymax=110,
        xtick={0, 10, 20, 30, 40, 50},
        ytick={0, 20, 40, 60, 80, 100},
        grid=major,
        grid style={gray!30},
        thick,
        legend pos=south west, legend style={font=\footnotesize},
]
\addplot[color=red, mark=triangle*, mark size=2, line width=1.2pt] coordinates {
        (0, 98.2) (5, 98.3) (10, 98.3) (15, 98.3)
        (20, 98.0) (25, 97.3) (30, 95.6) (35, 91.5)
        (40, 81.9) (45, 63.5) (50, 37.3)
};
\addlegendentry{Baseline}
\addplot[color=blue, mark=square*, mark size=2, line width=1.2pt] coordinates {
        (0, 98.2) (5, 98.2) (10, 97.9) (15, 97.4)
        (20, 96.4) (25, 94.4) (30, 91.0) (35, 83.9)
        (40, 70.7) (45, 49.1) (50, 33.3)
};
\addlegendentry{$\alpha{=}0.85,\;\beta{=}0.8$}
\addplot[color=gray, mark=diamond*, mark size=2, line width=1.2pt, dashed] coordinates {
        (0, 98.2) (5, 97.8) (10, 96.3) (15, 93.6)
        (20, 89.1) (25, 80.6) (30, 69.7) (35, 54.1)
        (40, 39.3) (45, 27.2) (50, 19.5)
};
\addlegendentry{$\alpha{=}0.5,\;\beta{=}0.5$}
\end{axis}
\end{tikzpicture}
\caption{Sufficiency evaluation on ResNet-CIFAR. Target probability as neurons outside the slice are progressively scaled down.}
\label{fig:counterfactual_sufficiency}
\end{figure}

\vspace{0.2cm}

Both evaluation methods confirm that our slices capture the relevant subnetwork for each class.
Pruning shows that the model retains classification accuracy when only the slice components are
kept.
Counterfactual reasoning shows that the identified neurons are necessary for the prediction and
that removing non-slice neurons has little effect.
The combined configuration ($\alpha{=}0.85$, $\beta{=}0.8$) performs nearly as well as the
unfiltered baseline in both evaluations while producing smaller slices.
The aggressive configuration degrades results in both cases, confirming that the thresholds should
not be set too low.

Having established that our slices are accurate, we now ask whether slices also reveal
interpretable patterns at the block and channel level.

\subsection*{RQ3: Do slicing results reveal interpretable class-specific patterns at channel and block level?}

Analyzing individual neuron contributions in a network with millions of neurons is neither
practical nor insightful.
We therefore investigate whether interpretable patterns emerge when aggregating contributions to
the block and channel level.
By interpretable patterns, we mean whether slices of the same class share similar block and channel
contributions, and whether semantically related classes show more overlap than unrelated ones.

Intuitively, semantically related classes should activate similar parts of the network.
For example, animal classes such as cat and dog share visual features like fur, legs, and eyes,
while vehicle classes such as automobile and truck share features like wheels and metallic
surfaces.
If such patterns are visible in the slice contributions, it would confirm that slices capture
meaningful structure and support the idea behind our hierarchical slicing approach.

To answer RQ3, we compute slices without our extension of block and channel slicing.
Instead, we aggregate the raw neuron contributions into their respective channels and blocks.
We use our custom ResNet model trained on CIFAR-10, as the dataset provides ten classes that split
naturally into animals and vehicles.
For each class, we slice 20 randomly sampled images that the model classifies correctly, resulting
in 200 slices in total.
These slices are then compared pairwise, first at the block level and then at the channel level.

\paragraph{Block-Level Interpretation.}
To compare slices at the block level, we focus exclusively on convolutional layers.
These layers contain the learned kernel weights that we consider as the actual neurons of the
network.
Each residual block $b \in \mathcal{B}$ contains two convolutional layers, each with a set of
channels $\mathcal{C}$.
For each channel $c \in \mathcal{C}$, we aggregate the contributions of its neurons
$\mathcal{N}_{\text{ch}}$ into a single value.
These channel contributions are then concatenated across both layers into a single vector per
block.
Since ResNet-CIFAR consists of 27 residual blocks, this produces 27 contribution vectors per image.

Now that we have a contribution vector for each block, the question is how to compare them across
different images and classes.
Since we are interested in whether two slices activate the same channels rather than whether their
absolute contribution values match, we use cosine similarity to compare the block vectors.
It measures the directional alignment between two vectors, independent of their magnitude:

\[
\text{cos}(\mathbf{a}, \mathbf{b}) = \frac{\mathbf{a} \cdot \mathbf{b}}{\|\mathbf{a}\| \, \|\mathbf{b}\|}
\]

In general, cosine similarity yields a value between $-1$ and $1$.
However, since we use absolute contribution values, all entries are non-negative and the similarity
ranges from $0$ to $1$:

\begin{itemize}
    \item $1$: the two vectors point in the same direction. Every channel has the same contribution value in both vectors.
    \item $0$: the two vectors are orthogonal. The channels that contribute to one slice do not contribute to the other.
\end{itemize}

We chose these four scenarios to test different levels of semantic similarity:
\begin{itemize}
    \item \textbf{Same class:} For each class in CIFAR-10, we compare slices of images belonging to the same class, e.g. for two different cat images. We expect their contribution vectors to be highly aligned.
    \item \textbf{Cat -- Dog:} As representatives of the animal classes, we compare slices of cat and dog images. Both share biological features such as fur and body shape. We expect high similarity, though not as high as the comparison of same class images.
    \item \textbf{Truck -- Automobile:} As representatives of the vehicle classes, slices of truck and automobile images are compared. Both share structural features like wheels and metallic surfaces. As with cat and dog, we expect their similarity to be high but below same-class level.
    \item \textbf{Cat -- Truck:} As a negative example, we pair one class from the animal group with one from the vehicle group. These classes have few visual features in common, and we expect the least similarity among our scenarios.
\end{itemize}

The block-level comparison results for ResNet-CIFAR are presented in
Figure~\ref{fig:similarity_by_block}, showing blocks 19 through 27. Blocks 1 through 24 show
near-zero cosine similarity across all four scenarios, indicating no class-specific structure at
the block level. From block 25 onward, the scenarios begin to diverge, with the clearest
differences appearing in block 27. Same-class pairs reach the highest similarity at 0.45,
suggesting that images of the same class activate similar channels in the final block. The
semantically related pairs cat--dog (0.15) and truck--automobile (0.11) also show notable
similarity, indicating that classes sharing visual features rely on partially overlapping channels.
As expected, cat--truck shows the lowest similarity at 0.03, as animals and vehicles share few
visual features.
While the absolute similarity values are relatively low, this is expected given that each block
contains 64 channels. Even if two classes share a few channels, the majority of channels differ,
which lowers the overall cosine similarity. The key finding is the clear ordering across scenarios,
which confirms that class-specific structure exists at the block level.

\begin{figure}[ht]
\centering
\begin{tikzpicture}
\begin{axis}[
        width=0.85\textwidth,
        height=0.50\textwidth,
        xlabel={Block},
        ylabel={Cosine Similarity},
        xmin=18.5, xmax=27.5,
        ymin=-0.05, ymax=0.5,
        xtick={19,20,21,22,23,24,25,26,27},
        grid=major,
        grid style={gray!30},
        thick,
        legend pos=north west, legend style={font=\footnotesize},
]
\addplot[color=green!60!black, mark=o, mark size=2, line width=1.2pt] coordinates {
        (19, 0.0004) (20, -0.0003) (21, 0.0003) (22, -0.0000)
        (23, 0.0004) (24, 0.0028) (25, 0.0158) (26, 0.0951) (27, 0.4486)
};
\addlegendentry{Same class}
\addplot[color=red, mark=square*, mark size=2, line width=1.2pt] coordinates {
        (19, -0.0004) (20, 0.0009) (21, 0.0005) (22, -0.0008)
        (23, 0.0006) (24, 0.0008) (25, 0.0107) (26, 0.0429) (27, 0.1480)
};
\addlegendentry{Cat -- Dog}
\addplot[color=blue, mark=triangle*, mark size=2, line width=1.2pt] coordinates {
        (19, -0.0006) (20, 0.0003) (21, 0.0004) (22, 0.0011)
        (23, -0.0008) (24, 0.0013) (25, 0.0069) (26, 0.0508) (27, 0.1085)
};
\addlegendentry{Truck -- Automobile}
\addplot[color=gray, mark=diamond*, mark size=2, line width=1.2pt] coordinates {
        (19, -0.0005) (20, 0.0013) (21, -0.0000) (22, -0.0003)
        (23, -0.0008) (24, -0.0009) (25, -0.0000) (26, 0.0144) (27, 0.0347)
};
\addlegendentry{Cat -- Truck}
\end{axis}
\end{tikzpicture}
\caption{Cosine similarity between slice contribution vectors for ResNet-CIFAR across blocks 19--27. Same-class pairs show the highest similarity in the final blocks, followed by semantically related pairs (cat--dog, truck--automobile). Cat--truck shows the lowest similarity throughout.}
\label{fig:similarity_by_block}
\end{figure}

These results align with the well-known observation that early layers in deep networks learn
generic, low-level features such as edges and textures, while later layers develop increasingly
class-specific representations~\cite{lecun2015deep, zeiler2014visualizing,
yosinski2014transferable}.
At the block level, class-specific patterns only become visible in the final blocks. This is where
the network begins to distinguish between semantic categories.
Overlapping patterns may exist in earlier layers as well. However, they may not be detectable at
the block level because contributions from many channels are aggregated into a single vector.

To investigate which individual channels drive these differences, we now zoom into the channel
level of the last block.


\paragraph{Channel-Level Interpretation.}
We now examine whether individual channels within block~27 show class-specific patterns.
Similar to the block-level interpretation, we consider the same semantically related pairs.
In addition, we extend the comparison to animal classes versus vehicle classes.

Block~27 contains two convolutional layers with 64 channels each. We analyzed both layers and found
that they show similar patterns. We therefore focus on the first convolutional layer as an example.
For each channel $c$ in the first convolutional layer, we average the absolute contribution values
across all 20 images of a class. This gives a single value per class per channel, representing how
strongly that channel contributes to slices of that class on average.
For each scenario, we select the two channels with the highest overlap between the two target
classes.

When comparing cat and dog slices, channels 23 and 31 emerge as the most shared channels
(Figure~\ref{fig:channel_cat_dog}).
Channel 23 shows a mean contribution of 0.31 for cat and 0.29 for dog, averaged across 20 images
per class. Channel 31 reaches 0.33 for cat and 0.38 for dog. While these values may appear low,
most channels show near-zero contributions for any given class, making these among the highest in
the layer.
However, other animal classes such as frog, horse, and deer also show notable contributions,
suggesting that these channels respond to shared animal features rather than being exclusive to cat
and dog.

\begin{figure}[ht]
\centering
\begin{tikzpicture}
\begin{axis}[
        ybar stacked, width=0.48\textwidth, height=5.5cm,
        bar width=8pt, ymin=0, ymax=0.85,
        ylabel={Contribution},
        title={Channel 23},
        symbolic x coords={frog,cat,dog,deer,horse,airplane,bird,auto,truck,ship},
        xtick=data, xticklabel style={rotate=45, anchor=east, font=\footnotesize},
        ymajorgrids=true, grid style={gray!30}, legend entries={},
]
\addplot[fill=red!70, draw=black, line width=0.3pt] coordinates {
        (frog,0) (cat,0.3117) (dog,0.2891) (deer,0) (horse,0)
        (airplane,0) (bird,0) (auto,0) (truck,0) (ship,0)};
\addplot[fill=gray!40, draw=black, line width=0.3pt] coordinates {
        (frog,0.4352) (cat,0) (dog,0) (deer,0.0719) (horse,0.0555)
        (airplane,0.0414) (bird,0.0406) (auto,0.0320) (truck,0.0102) (ship,0.0102)};
\end{axis}
\end{tikzpicture}%
\hfill%
\begin{tikzpicture}
\begin{axis}[
        ybar stacked, width=0.48\textwidth, height=5.5cm,
        bar width=8pt, ymin=0, ymax=0.85,
        title={Channel 31},
        symbolic x coords={dog,horse,cat,frog,deer,bird,airplane,auto,truck,ship},
        xtick=data, xticklabel style={rotate=45, anchor=east, font=\footnotesize},
        ymajorgrids=true, grid style={gray!30}, legend entries={},
]
\addplot[fill=red!70, draw=black, line width=0.3pt] coordinates {
        (dog,0.3820) (horse,0) (cat,0.3336) (frog,0) (deer,0)
        (bird,0) (airplane,0) (auto,0) (truck,0) (ship,0)};
\addplot[fill=gray!40, draw=black, line width=0.3pt] coordinates {
        (dog,0) (horse,0.3508) (cat,0) (frog,0.1656) (deer,0.1414)
        (bird,0.0945) (airplane,0.0523) (auto,0.0375) (truck,0.0289) (ship,0.0258)};
\end{axis}
\end{tikzpicture}
\caption{Channels 23 and 31: the two channels with the largest overlap between cat and dog in block~27 of ResNet-CIFAR. Cat and dog show high contributions, though other animal classes also contribute notably.}
\label{fig:channel_cat_dog}
\end{figure}

The same analysis for truck and automobile is shown in Figure~\ref{fig:channel_truck_auto}.
Channels 50 and 4, different from the cat--dog channels, respond strongly to truck and automobile
while all other classes remain low. 
These channels are exclusively high for truck and automobile, suggesting that they capture features
specific to these two classes that are not shared by other classes.

\begin{figure}[ht]
\centering
\begin{tikzpicture}
\begin{axis}[
        ybar stacked, width=0.48\textwidth, height=5.5cm,
        bar width=8pt, ymin=0, ymax=0.85,
        ylabel={Contribution},
        title={Channel 50},
        symbolic x coords={auto,truck,ship,bird,airplane,horse,dog,cat,frog,deer},
        xtick=data, xticklabel style={rotate=45, anchor=east, font=\footnotesize},
        ymajorgrids=true, grid style={gray!30}, legend entries={},
]
\addplot[fill=blue!60, draw=black, line width=0.3pt] coordinates {
        (auto,0.3297) (truck,0.3047) (ship,0) (bird,0) (airplane,0)
        (horse,0) (dog,0) (cat,0) (frog,0) (deer,0)};
\addplot[fill=gray!40, draw=black, line width=0.3pt] coordinates {
        (auto,0) (truck,0) (ship,0.0594) (bird,0.0391) (airplane,0.0289)
        (horse,0.0273) (dog,0.0250) (cat,0.0234) (frog,0.0227) (deer,0.0211)};
\end{axis}
\end{tikzpicture}%
\hfill%
\begin{tikzpicture}
\begin{axis}[
        ybar stacked, width=0.48\textwidth, height=5.5cm,
        bar width=8pt, ymin=0, ymax=0.85,
        title={Channel 4},
        symbolic x coords={auto,truck,cat,airplane,horse,ship,bird,dog,deer,frog},
        xtick=data, xticklabel style={rotate=45, anchor=east, font=\footnotesize},
        ymajorgrids=true, grid style={gray!30}, legend entries={},
]
\addplot[fill=blue!60, draw=black, line width=0.3pt] coordinates {
        (auto,0.4211) (truck,0.4039) (cat,0) (airplane,0) (horse,0)
        (ship,0) (bird,0) (dog,0) (deer,0) (frog,0)};
\addplot[fill=gray!40, draw=black, line width=0.3pt] coordinates {
        (auto,0) (truck,0) (cat,0.1273) (airplane,0.1258) (horse,0.1055)
        (ship,0.0789) (bird,0.0781) (dog,0.0742) (deer,0.0727) (frog,0.0453)};
\end{axis}
\end{tikzpicture}
\caption{Channels 50 and 4: the two channels with the largest overlap between truck and automobile in block~27 of ResNet-CIFAR. Truck and automobile show high contributions while other classes remain low.}
\label{fig:channel_truck_auto}
\end{figure}

The pairwise comparisons show that semantically related classes share channels at the channel
level. 
For truck and automobile, this sharing is exclusive, while cat and dog channels also respond to
other animal classes. 
This raises the question whether channels specialize not just on class pairs, but on broader
semantic groups. 
We therefore extend the analysis and compare animal classes against vehicle classes. 
The two channels with the largest overlap across animal classes are shown in
Figure~\ref{fig:channel_animals}.


\begin{figure}[ht]
\centering
\begin{tikzpicture}
\begin{axis}[
        ybar stacked, width=0.48\textwidth, height=5.5cm,
        bar width=8pt, ymin=0, ymax=0.85,
        ylabel={Contribution},
        title={Channel 49},
        symbolic x coords={dog,frog,deer,cat,horse,bird,airplane,auto,ship,truck},
        xtick=data, xticklabel style={rotate=45, anchor=east, font=\footnotesize},
        ymajorgrids=true, grid style={gray!30}, legend entries={},
]
\addplot[fill=red!70, draw=black, line width=0.3pt] coordinates {
        (dog,0.4602) (frog,0.4297) (deer,0.3844) (cat,0.3836) (horse,0.3187)
        (bird,0.3148) (airplane,0) (auto,0) (ship,0) (truck,0)};
\addplot[fill=blue!60, draw=black, line width=0.3pt] coordinates {
        (dog,0) (frog,0) (deer,0) (cat,0) (horse,0)
        (bird,0) (airplane,0.1531) (auto,0.1031) (ship,0.0883) (truck,0.0773)};
\end{axis}
\end{tikzpicture}%
\hfill%
\begin{tikzpicture}
\begin{axis}[
        ybar stacked, width=0.48\textwidth, height=5.5cm,
        bar width=8pt, ymin=0, ymax=0.85,
        title={Channel 16},
        symbolic x coords={frog,deer,horse,cat,dog,bird,airplane,truck,auto,ship},
        xtick=data, xticklabel style={rotate=45, anchor=east, font=\footnotesize},
        ymajorgrids=true, grid style={gray!30}, legend entries={},
]
\addplot[fill=red!70, draw=black, line width=0.3pt] coordinates {
        (frog,0.3602) (deer,0.3281) (horse,0.2875) (cat,0.2102) (dog,0.2062)
        (bird,0.1016) (airplane,0) (truck,0) (auto,0) (ship,0)};
\addplot[fill=blue!60, draw=black, line width=0.3pt] coordinates {
        (frog,0) (deer,0) (horse,0) (cat,0) (dog,0)
        (bird,0) (airplane,0.0984) (truck,0.0961) (auto,0.0391) (ship,0.0391)};
\end{axis}
\end{tikzpicture}
\caption{Channels 49 and 16: the two channels with the largest overlap across animal classes in block~27 of ResNet-CIFAR. All six animal classes show high contributions while vehicle classes remain low.}
\label{fig:channel_animals}
\end{figure}

Channels 49 and 16 both show high contributions for all six animal classes while vehicle classes
remain low. This suggests that these channels respond to features common to animals.

On the other hand, Figure~\ref{fig:channel_vehicles} shows the two channels with the largest
overlap across vehicle classes.
Channels 11 and 28 show the opposite pattern: all four vehicle classes contribute strongly while
animal classes remain low. This suggests that these channels respond to features common to
vehicles.

\begin{figure}[ht]
\centering
\begin{tikzpicture}
\begin{axis}[
        ybar stacked, width=0.48\textwidth, height=6cm,
        bar width=8pt, ymin=0, ymax=0.85,
        ylabel={Contribution},
        title={Channel 11},
        symbolic x coords={truck,ship,airplane,auto,cat,horse,bird,deer,dog,frog},
        xtick=data, xticklabel style={rotate=45, anchor=east, font=\footnotesize},
        ymajorgrids=true, grid style={gray!30}, legend entries={},
]
\addplot[fill=blue!60, draw=black, line width=0.3pt] coordinates {
        (truck,0.5766) (ship,0.5617) (airplane,0.4953) (auto,0.4336)
        (cat,0) (horse,0) (bird,0) (deer,0) (dog,0) (frog,0)};
\addplot[fill=red!70, draw=black, line width=0.3pt] coordinates {
        (truck,0) (ship,0) (airplane,0) (auto,0)
        (cat,0.1156) (horse,0.0883) (bird,0.0844) (deer,0.0625) (dog,0.0484) (frog,0.0125)};
\end{axis}
\end{tikzpicture}%
\hfill%
\begin{tikzpicture}
\begin{axis}[
        ybar stacked, width=0.48\textwidth, height=6cm,
        bar width=8pt, ymin=0, ymax=0.85,
        title={Channel 28},
        symbolic x coords={ship,airplane,auto,frog,truck,horse,deer,cat,bird,dog},
        xtick=data, xticklabel style={rotate=45, anchor=east, font=\footnotesize},
        ymajorgrids=true, grid style={gray!30}, legend entries={},
]
\addplot[fill=blue!60, draw=black, line width=0.3pt] coordinates {
        (ship,0.4703) (airplane,0.4336) (auto,0.3711) (frog,0)
        (truck,0.1883) (horse,0) (deer,0) (cat,0) (bird,0) (dog,0)};
\addplot[fill=red!70, draw=black, line width=0.3pt] coordinates {
        (ship,0) (airplane,0) (auto,0) (frog,0.1984)
        (truck,0) (horse,0.1789) (deer,0.1273) (cat,0.1039) (bird,0.0711) (dog,0.0375)};
\end{axis}
\end{tikzpicture}
\caption{Channels 11 and 28: the two channels with the largest overlap across vehicle classes in block~27 of ResNet-CIFAR. All four vehicle classes show high contributions while animal classes remain low.}
\label{fig:channel_vehicles}
\end{figure}

In summary, when aggregating slice contributions to the block and channel level, interpretable
patterns emerge. At the block level, the final blocks show differentiation between semantically
related and unrelated classes. Zooming into the channel level reveals that individual channels are
shared between similar classes or even entire semantic groups such as animals or vehicles. We
demonstrated this exemplarily on ResNet-CIFAR using representative class pairs from the animal and
vehicle groups. These patterns suggest that slices capture meaningful class-specific structure,
which could inform decisions such as which channels or blocks to preserve during filtering or
pruning.

\vspace{0.2cm}

Taking the results of all three research questions together, we evaluated both the efficiency and
quality of our approach as well as the interpretability of the produced slices.
RQ1 shows that our vectorized implementation combined with block and channel slicing significantly
reduces runtime.
RQ2 confirms that the resulting slices retain relevant information, demonstrated through pruning
and counterfactual experiments.
RQ3 demonstrates that the slices reveal interpretable, class-specific patterns at the last blocks
and within their channels.
We conclude the evaluation by examining potential threats to the validity of these results.

\section{Threats to Validity}

This section discusses factors that may affect the validity of our results.
We structure the discussion into internal, external, and construct validity.
Internal validity concerns whether our experimental setup could introduce bias or confounding
factors.
External validity addresses whether our findings generalize beyond the specific models and datasets
used.
Construct validity examines our evaluation methods and design choices.

\paragraph{Internal Validity.}
Internal validity threats stem from the parameter choices in our experiments.
Three parameters influence the results: the contribution threshold $\theta$, the block threshold
$\beta$, and the channel threshold $\alpha$.
Different combinations lead to different slice sizes and quality outcomes.
We addressed this by testing a broad range of values for $\alpha$ and $\beta$ and reporting results
across multiple configurations.
For $\theta$, we adopted the value of $0.2$ based on the original NNSlicer implementation.
How a different value of $\theta$ interacts with $\alpha$ and $\beta$ remains to be tested.

Another point concerns the sample sizes used in the evaluation.
In RQ2, we compute slices from 10 images per class, and in RQ3 from 20 images per class.
Larger samples could lead to more robust findings.

\paragraph{External Validity.}
External validity threats concern the generalizability of our approach to other architectures and
datasets.
Block-based slicing works for any ResNet regardless of depth or size, which we demonstrated by
evaluating three models of different depths in RQ1.
The vectorized backward analysis and channel slicing also apply to non-residual convolutional
networks such as LeNet, since they only require convolutional layers.
Other architectures such as DenseNets or ConvNeXt also organize layers into blocks. However,
adapting block slicing to architectures without skip connections remains an open direction.

A separate concern is the generalizability of the slice quality evaluation.
We evaluated slice quality on a single custom ResNet with 27 residual blocks trained on CIFAR-10,
which we chose to ensure reproducibility. We did not vary the number of blocks or the model depth
for these experiments.
Whether similar results hold for models with fewer or more blocks, other datasets, or larger
architectures is open for future work. 
Evaluating on larger-scale datasets and models would likely require stronger hardware with GPU
support to compute slices over a larger number of images.

\paragraph{Construct Validity.}
There is no ground truth for what constitutes a correct neural network slice.
We therefore evaluate slice quality indirectly through two methods.
Model pruning, which was also used in the original NNSlicer evaluation, tests whether a pruned
model still achieves high accuracy on the target class. If so, the slice must have retained the
most relevant components.
As a second method, we apply counterfactual reasoning, which has not previously been used to
evaluate neural network slices. It complements pruning by testing how the prediction changes when
parts of the slice are progressively modified.
Both methods assess whether the slice retains information that is important for classification, but
they do not guarantee that the slice captures all and only the relevant neurons.

Our slicing approach also relies on specific design choices about how we define relevance.
We treat the aggregated delta as a relevance score for blocks and channels: components with a delta
close to zero are considered unimportant for the current input and may be filtered out.
However, a low delta only indicates that activations are close to their profiling mean and does not
rule out that a component consistently supports correct predictions.
Similarly, channel slicing removes entire channels at once, which means that individual important
neurons within a low-delta channel could be lost.


\chapter{Summary and Future Work}

This chapter summarizes the contributions and results of this work and outlines directions for
future research.

\section{Summary}

In this thesis, we introduced a block-based extension of the NNSlicer dynamic slicing approach.
We analyzed the model at higher levels of granularity and identified blocks and channels that are
less relevant for a specific target class.
To measure block and channel relevance, we aggregate neuron deltas, where each delta measures how
much the activation of a neuron deviates from its mean.
These deltas are aggregated to the channel and block level.
Based on these aggregated deltas, we introduced two filtering mechanisms with separate
hyperparameters.
Blocks whose aggregated delta falls below a threshold~$\beta$ are skipped entirely, and the
contribution flow is routed through the skip connections.
For the remaining blocks, channel slicing filters channels whose aggregated delta falls below a
threshold~$\alpha$.

For the neuron-level analysis on the remaining network, we proposed a vectorized reimplementation
of NNSlicer.
Instead of computing neuron contributions individually in loops, our implementation uses tensor
operations in PyTorch.

We evaluated our approach along three research questions addressing runtime efficiency, slice
quality, and interpretability.
Runtime was evaluated on ResNet-18, ResNet-34, and ResNet-101 pretrained on ImageNet.
For slice quality and interpretability we trained a custom ResNet on CIFAR-10.

The vectorized reimplementation alone reduces slicing time from over 540 seconds to approximately
two seconds for ResNet-18.
Even large models like ResNet-101 with 44.5M parameters can now be sliced in about 7.7 seconds.
Block slicing further reduces runtime and slice size by skipping entire blocks. At $\beta = 0.80$,
ResNet-101 drops from 7.67\,s to 6.50\,s while the neuron ratio decreases from 41.8\,\% to
35.8\,\%.
Channel slicing does not reduce runtime, as the tensor dimensions remain unchanged, but it further
reduces slice size by zeroing out low-energy channels.
Combining both mechanisms with $\alpha = 0.85$ and $\beta = 0.80$ produces slices that are both
smaller and faster to compute than the baseline.
Both block and channel thresholds can be adjusted to produce more aggressive slices. Lower values
for $\beta$ and $\alpha$ yield smaller and faster slices at the cost of potentially losing relevant
components.

To validate that the filtered slices still capture class-relevant information, we applied two
evaluation strategies.
First, we pruned the model by removing all parameters outside the slice and evaluated the accuracy
of the pruned model.
Block pruning at $\beta = 0.80$ reduces the model to 39.7\,\% of its parameters while target
accuracy remains at 91.9\,\%, close to the baseline of 93.6\,\%.
Channel pruning alone degraded more quickly, with target accuracy dropping to 70.5\,\% at $\alpha =
0.80$.
Second, we applied counterfactual reasoning to ask what happens if parts of the slice are altered.
We progressively scaled down neurons inside the slice (necessity) and outside the slice
(sufficiency) to measure how the prediction changes.
The combined configuration ($\alpha{=}0.85$, $\beta{=}0.8$) performed nearly as well as the
unfiltered baseline in both tests, while an aggressive configuration ($\alpha{=}0.5$,
$\beta{=}0.5$) degraded in both cases.
Both strategies confirm that conservative thresholds preserve high slice quality.

Finally, we investigated whether slices reveal interpretable class-specific patterns.
For this analysis, we used the vectorized neuron-level approach without block or channel filtering.
At the block level, we compared contribution vectors across class pairs using cosine similarity and
found that similarity between classes only emerges in the final layers, consistent with the known
progression from generic to class-specific features in deep networks.
We then examined individual channels in the last residual block for class pairs where we expected
high overlap, such as cat--dog and truck--automobile.
Certain channels were shared between semantically related classes, suggesting specialization for
broader categories such as animals and vehicles.

Overall, our results show that vectorized computation and block-based filtering make dynamic
slicing practical for modern ResNet architectures.
The resulting slices are fast to compute, retain classification-relevant information, and provide
first evidence of interpretable structure at block and channel level.


\section{Future Work}

Building on our results, we identify several directions for future work.

The original NNSlicer authors propose processing multiple slices in parallel to reduce runtime.
Our vectorized implementation already outperforms their batched approach for individual slices, but
combining both strategies could reduce the per-slice cost even further when many slices need to be
computed simultaneously.

Another direction is the extension to other architectures. The block-based approach is currently
limited to residual networks.
However, other network families contain structural units that could serve as slicing boundaries.
ConvNeXt uses residual blocks with skip connections similar to ResNets~\cite{liu2022convnet}.
Transformers also use residual connections around their attention and feed-forward sublayers,
though the slicing procedure would need to be adapted to handle attention mechanisms instead of
convolutions~\cite{vaswani2017attention}.
DenseNets connect each layer to all preceding layers, which prevents skipping individual blocks but
could still allow filtering at the layer or channel level~\cite{huang2017densely}.

A third direction concerns the scale of the evaluation. The slice quality and interpretability
experiments were conducted on CIFAR-10 using a single custom ResNet.
Evaluating pruning and counterfactual experiments on larger-scale models and datasets such as
ImageNet would provide stronger evidence that the approach generalizes.
Similarly, applying the interpretability analysis to datasets with more classes and finer semantic
distinctions could reveal more patterns in block and channel comparisons.
Additionally, varying the number of residual blocks beyond the 27 used in our custom model would
help assess how block depth affects slice quality and interpretability.

On the parameter side, $\theta$ is fixed at $0.2$ throughout all our experiments. A systematic
study of how $\theta$ interacts with the block and channel thresholds $\alpha$ and $\beta$ could
provide guidance on how to choose these parameters jointly.

Finally, the runtime improvements achieved in this work extend the applicability of neural network
slicing to larger models.
Zhang et al. already demonstrated that slices can be used for adversarial input detection, where
slices of adversarial samples differ from those of normal samples~\cite{zhang2020dynamic} .
However, this was only feasible for small networks due to the high runtime.
With slicing times of a few seconds, such applications can now scale to larger networks.

\backmatter
\printbibliography

\end{document}
